% Generated by the blueprint tool (blueprint-pdf/1) from a blueprint snapshot; edit the
% Markdown sources, not this file.  Compile with `latexmk -pdf` or by running
% `pdflatex` twice (the second run settles the contents and the page references).
% Images, if any, are read from assets/ next to this file.
\documentclass[11pt]{article}
\usepackage[T1]{fontenc}
\usepackage[utf8]{inputenc}
\usepackage{lmodern}
\usepackage{textcomp}
\usepackage[a4paper,margin=2.5cm,headheight=14pt]{geometry}
\usepackage{amsmath,amssymb,amsthm}
\usepackage{amscd}  % KaTeX's CD diagrams; part of amsmath, so always there
\usepackage{mathtools}
\usepackage{bm}
\usepackage{graphicx}
\usepackage{xcolor}
\usepackage{array,booktabs,tabularx}
\usepackage{fancyvrb}
\usepackage[normalem]{ulem}
\usepackage{microtype}
% Optional, and only with their fonts: a small TeX installation has the
% packages without the fonts, which fails at the first formula.
\IfFileExists{mathrsfs.sty}{\IfFileExists{rsfs10.tfm}{\usepackage{mathrsfs}}{}}{}
\IfFileExists{stmaryrd.sty}{\IfFileExists{stmary10.tfm}{\usepackage{stmaryrd}}{}}{}
\IfFileExists{cancel.sty}{\usepackage{cancel}}{}
\usepackage[hyphens]{url}
\usepackage[pdfusetitle,bookmarksnumbered=false]{hyperref}
\usepackage{bookmark}
\definecolor{bplink}{HTML}{1F4E9C}
\definecolor{bpmuted}{HTML}{5F676E}
\hypersetup{colorlinks=true,linkcolor=bplink,urlcolor=bplink,citecolor=bplink}
\setlength{\parindent}{0pt}
\setlength{\parskip}{0.55em plus 0.2em minus 0.1em}
% Lists as tight as the page's.
\makeatletter
\def\@listi{\leftmargin\leftmargini\parsep 0pt\topsep 0.3em plus 0.1em\itemsep 0.1em plus 0.1em}
\let\@listI\@listi
\def\@listii{\leftmargin\leftmarginii\labelwidth\leftmarginii\advance\labelwidth-\labelsep\parsep 0pt\topsep 0.1em\itemsep 0.1em}
% An ordered list that does not start at 1.
\newcommand{\bpstartat}[1]{\setcounter{\@enumctr}{\numexpr#1-1\relax}}
\makeatother
\setcounter{tocdepth}{3}
\emergencystretch=3em
\allowdisplaybreaks

% A numbered statement, set as a paper sets one: "Theorem 2.1 (Title)."
% #1 is the lead ("Theorem~2.1"), #2 what follows it in roman.
\newenvironment{bpstatement}[2]{\par\addvspace{0.6\baselineskip}\noindent\textbf{#1}#2\textbf{.}\enspace\ignorespaces}
  {\par\addvspace{0.4\baselineskip}}
\newenvironment{bpstatementit}[2]{\begin{bpstatement}{#1}{#2}\itshape\ignorespaces}{\end{bpstatement}}
% A step: a relation with prose of its own, read under the object it starts from.
\newenvironment{bpstep}[1]{\par\addvspace{0.3\baselineskip}\begin{list}{}{\leftmargin=1.6em\rightmargin=0pt\topsep=0pt}\item\relax
  {\small\color{bpmuted}#1}\par\nopagebreak}{\end{list}\par}
% Notes about where an object is written out.
\newenvironment{bpnote}{\par\small\color{bpmuted}\itshape}{\par}
\newcommand{\bpomitted}[1]{\par{\small\color{bpmuted}\itshape #1\par}}
\newcommand{\bpbodyheading}[1]{\par\addvspace{0.4\baselineskip}\noindent\textbf{#1}\par\nopagebreak}
\newcommand{\bprule}{\par\medskip\noindent\hfil\rule{0.3\linewidth}{0.4pt}\hfil\par\medskip}
\newcommand{\bpboundary}[1]{\par\nopagebreak{\small\color{bpmuted}#1\par}\nopagebreak}
% A reference to something this export leaves out, and one that resolves nowhere.
\newcommand{\bpout}[1]{#1\textsuperscript{\hyperlink{bp-outside}{\ensuremath{\ast}}}}
\newcommand{\bpbroken}[1]{\textcolor{red!70!black}{[#1?]}}
\newcommand{\bpslug}[1]{\texttt{#1}}
% A bibliography entry: label, anchor, entry; the label hangs in the margin.
\newcommand{\bpbibitem}[3]{\par\setbox0\hbox{[#1]\enspace}\ifdim\wd0<5em\setbox0\hbox to 5em{[#1]\hfil}\fi
  \hangindent=\wd0 \hangafter=1 \noindent\box0\phantomsection\label{#2}#3\par\smallskip}
% A character the export knows no way to set.
\newcommand{\bpglyph}[1]{\fbox{\scriptsize #1}}
% Images, at their own size unless they are wider than the line.
\newsavebox{\bpimagebox}
\newcommand{\bpimage}[1]{\sbox{\bpimagebox}{\includegraphics{#1}}%
  \ifdim\wd\bpimagebox>\linewidth\resizebox{\linewidth}{!}{\usebox{\bpimagebox}}\else\usebox{\bpimagebox}\fi}
\newcommand{\bpmissingimage}[1]{\fbox{\parbox{0.8\linewidth}{\centering\small\color{bpmuted}Image not included: #1}}}
\DefineVerbatimEnvironment{bpverbatim}{Verbatim}{fontsize=\small,xleftmargin=1em}
\renewcommand{\qedsymbol}{\ensuremath{\blacksquare}}
\hypersetup{pdftitle={Weil Conjectures — 1 Introduction},pdfsubject={Section 1 Introduction; all levels},pdfcreator={blueprint}}
% KaTeX commands the website knows
\providecommand{\R}{\mathbb{R}}\providecommand{\N}{\mathbb{N}}\providecommand{\Z}{\mathbb{Z}}
\providecommand{\Q}{\mathbb{Q}}\providecommand{\C}{\mathbb{C}}
\providecommand{\Reals}{\mathbb{R}}\providecommand{\reals}{\mathbb{R}}\providecommand{\Complex}{\mathbb{C}}
\providecommand{\cnums}{\mathbb{C}}\providecommand{\natnums}{\mathbb{N}}\providecommand{\Bbb}[1]{\mathbb{#1}}
\providecommand{\bold}[1]{\mathbf{#1}}\providecommand{\frak}[1]{\mathfrak{#1}}\providecommand{\mathbbm}[1]{\mathbb{#1}}
\providecommand{\lang}{\langle}\providecommand{\rang}{\rangle}\providecommand{\bra}[1]{\langle #1\rvert}
\providecommand{\ket}[1]{\lvert #1\rangle}\providecommand{\braket}[1]{\langle #1\rangle}
\providecommand{\set}[1]{\{#1\}}\providecommand{\Set}[1]{\left\{#1\right\}}
\providecommand{\infin}{\infty}\providecommand{\isin}{\in}\providecommand{\plusmn}{\pm}\providecommand{\sub}{\subset}
\providecommand{\sube}{\subseteq}\providecommand{\supe}{\supseteq}\providecommand{\exist}{\exists}
\providecommand{\bull}{\bullet}\providecommand{\larr}{\leftarrow}\providecommand{\rarr}{\rightarrow}
\providecommand{\lrarr}{\leftrightarrow}\providecommand{\Larr}{\Leftarrow}\providecommand{\Rarr}{\Rightarrow}
\providecommand{\Lrarr}{\Leftrightarrow}\providecommand{\uarr}{\uparrow}\providecommand{\darr}{\downarrow}
\providecommand{\harr}{\leftrightarrow}\providecommand{\lArr}{\Leftarrow}\providecommand{\rArr}{\Rightarrow}
\providecommand{\hArr}{\Leftrightarrow}\providecommand{\alef}{\aleph}\providecommand{\alefsym}{\aleph}
\providecommand{\real}{\Re}\providecommand{\image}{\Im}\providecommand{\weierp}{\wp}
\providecommand{\Alpha}{\mathrm{A}}\providecommand{\Beta}{\mathrm{B}}\providecommand{\Epsilon}{\mathrm{E}}
\providecommand{\Zeta}{\mathrm{Z}}\providecommand{\Eta}{\mathrm{H}}\providecommand{\Iota}{\mathrm{I}}
\providecommand{\Kappa}{\mathrm{K}}\providecommand{\Mu}{\mathrm{M}}\providecommand{\Nu}{\mathrm{N}}
\providecommand{\Omicron}{\mathrm{O}}\providecommand{\Rho}{\mathrm{P}}\providecommand{\Tau}{\mathrm{T}}
\providecommand{\Chi}{\mathrm{X}}\providecommand{\omicron}{o}\providecommand{\gt}{>}\providecommand{\lt}{<}
\providecommand{\And}{\&}\providecommand{\hdashline}{\hline}\providecommand{\operatornamewithlimits}[1]{\operatorname*{#1}}
\providecommand{\llbracket}{[\![}\providecommand{\rrbracket}{]\!]}\providecommand{\cancel}[1]{#1}
\providecommand{\mathscr}{\mathcal}
% the characters beyond ASCII this document uses
\DeclareUnicodeCharacter{2160}{\textup{I}}
\DeclareUnicodeCharacter{2161}{\textup{II}}
\DeclareUnicodeCharacter{2162}{\textup{III}}
% the project's macros ([katex.macros] in blueprint.toml)
\providecommand{\A}{}\renewcommand{\A}{\mathbf A}
\providecommand{\AEn}{}\renewcommand{\AEn}{A_{E,n}}
\providecommand{\Aaf}{}\renewcommand{\Aaf}{\mathbf A}
\providecommand{\Ab}{}\renewcommand{\Ab}{\mathsf{Ab}}
\providecommand{\AffSch}{}\renewcommand{\AffSch}{\operatorname{AffSch}}
\providecommand{\Aone}{}\renewcommand{\Aone}{\mathbf A^1}
\providecommand{\Bl}{}\renewcommand{\Bl}{\operatorname{Bl}}
\providecommand{\Br}{}\renewcommand{\Br}{\mathrm{Br}}
\providecommand{\C}{}\renewcommand{\C}{\mathbb C}
\providecommand{\Cech}{}\renewcommand{\Cech}{\check{H}}
\providecommand{\Coext}{}\renewcommand{\Coext}{\operatorname{Coext}}
\providecommand{\Cone}{}\renewcommand{\Cone}{\operatorname{Cone}}
\providecommand{\Cons}{}\renewcommand{\Cons}{\operatorname{Cons}}
\providecommand{\Cref}{}\renewcommand{\Cref}[1]{\text{#1}}
\providecommand{\D}{}\renewcommand{\D}{D}
\providecommand{\DLAFT}{}\renewcommand{\DLAFT}{\operatorname{FT}}
\providecommand{\DLAFr}{}\renewcommand{\DLAFr}{\operatorname{Fr}}
\providecommand{\DLAQ}{}\renewcommand{\DLAQ}{{\overline{\mathbf Q}_\ell}}
\providecommand{\DLRQ}{}\renewcommand{\DLRQ}{{\overline{\mathbf Q}_\ell}}
\providecommand{\Dbc}{}\renewcommand{\Dbc}{D^b_c}
\providecommand{\Dc}{}\renewcommand{\Dc}{D_{\mathrm{cons}}}
\providecommand{\Dctf}{}\renewcommand{\Dctf}{D^b_{ctf}}
\providecommand{\Dic}{}\renewcommand{\Dic}{D_{\IndCons}}
\providecommand{\Dplus}{}\renewcommand{\Dplus}{D^{+}}
\providecommand{\End}{}\renewcommand{\End}{\operatorname{End}}
\providecommand{\Et}{}\renewcommand{\Et}{\mathrm{\acute{e}t}}
\providecommand{\Ex}{}\renewcommand{\Ex}{\mathrm{Ex}}
\providecommand{\Ext}{}\renewcommand{\Ext}{\operatorname{Ext}}
\providecommand{\F}{}\renewcommand{\F}{\mathbb F}
\providecommand{\FT}{}\renewcommand{\FT}{\operatorname{FT}}
\providecommand{\Fbar}{}\renewcommand{\Fbar}{{\overline{\mathbf F}_q}}
\providecommand{\Fell}{}\renewcommand{\Fell}{\mathbf F_\ell}
\providecommand{\Fib}{}\renewcommand{\Fib}{\operatorname{Fib}}
\providecommand{\Fix}{}\renewcommand{\Fix}{\operatorname{Fix}}
\providecommand{\Fp}{}\renewcommand{\Fp}{{\mathbf F_{p}}}
\providecommand{\Fq}{}\renewcommand{\Fq}{{\mathbf F_q}}
\providecommand{\Fqn}{}\renewcommand{\Fqn}{{\mathbf F_{q^{n}}}}
\providecommand{\Frob}{}\renewcommand{\Frob}{\operatorname{Frob}}
\providecommand{\GL}{}\renewcommand{\GL}{\operatorname{GL}}
\providecommand{\Gal}{}\renewcommand{\Gal}{\operatorname{Gal}}
\providecommand{\Hens}{}\renewcommand{\Hens}{\operatorname{Hens}}
\providecommand{\Hilb}{}\renewcommand{\Hilb}{\operatorname{Hilb}}
\providecommand{\Hom}{}\renewcommand{\Hom}{\operatorname{Hom}}
\providecommand{\IndCons}{}\renewcommand{\IndCons}{\operatorname{IndCons}}
\providecommand{\Kar}{}\renewcommand{\Kar}{\operatorname{Kar}}
\providecommand{\LT}{}\renewcommand{\LT}{\operatorname{LT}}
\providecommand{\Lie}{}\renewcommand{\Lie}{\operatorname{Lie}}
\providecommand{\Loc}{}\renewcommand{\Loc}{\operatorname{Loc}}
\providecommand{\Mod}{}\renewcommand{\Mod}{\mathsf{Mod}}
\providecommand{\N}{}\renewcommand{\N}{\mathbb{N}}
\providecommand{\Needspace}{}\renewcommand{\Needspace}{\par}
\providecommand{\O}{}\renewcommand{\O}{{\mathcal{O}}}
\providecommand{\Ob}{}\renewcommand{\Ob}{\mathrm{Ob}}
\providecommand{\Oe}{}\renewcommand{\Oe}{\mathcal O_E}
\providecommand{\PP}{}\renewcommand{\PP}{\mathbf{P}}
\providecommand{\Perf}{}\renewcommand{\Perf}{\operatorname{Perf}}
\providecommand{\Pic}{}\renewcommand{\Pic}{\mathrm{Pic}}
\providecommand{\Pm}{}\renewcommand{\Pm}{\mathbf P^{m}}
\providecommand{\Pone}{}\renewcommand{\Pone}{\mathbf P^{1}}
\providecommand{\Pqn}{}\renewcommand{\Pqn}{\mathbf P^{n}}
\providecommand{\ProFin}{}\renewcommand{\ProFin}{\operatorname{ProFin}}
\providecommand{\Proet}{}\renewcommand{\Proet}{\mathrm{pro\acute{e}t}}
\providecommand{\Q}{}\renewcommand{\Q}{\mathbb Q}
\providecommand{\Qb}{}\renewcommand{\Qb}{\overline{\mathbf Q}_{\ell}}
\providecommand{\Qbar}{}\renewcommand{\Qbar}{\overline{\mathbf Q}}
\providecommand{\QbarX}{}\renewcommand{\QbarX}{\overline{\mathbf Q}_{\ell,X}}
\providecommand{\Qbarl}{}\renewcommand{\Qbarl}{{\overline{\mathbf Q}_{\ell}}}
\providecommand{\Qell}{}\renewcommand{\Qell}{{\mathbf{Q}_{\ell}}}
\providecommand{\Ql}{}\renewcommand{\Ql}{\mathbf Q_{\ell}}
\providecommand{\Qlb}{}\renewcommand{\Qlb}{{\overline{\mathbf Q}_\ell}}
\providecommand{\Qlbar}{}\renewcommand{\Qlbar}{\overline{\mathbf Q}_\ell}
\providecommand{\R}{}\renewcommand{\R}{\mathbb{R}}
\providecommand{\RG}{}\renewcommand{\RG}{\operatorname{\mathrm{R}\Gamma}}
\providecommand{\RGamma}{}\renewcommand{\RGamma}{R\Gamma}
\providecommand{\RGammac}{}\renewcommand{\RGammac}{R\Gamma_c}
\providecommand{\RHom}{}\renewcommand{\RHom}{R\mathcal{H}om}
\providecommand{\RHomX}{}\renewcommand{\RHomX}{R\mathcal H om_X}
\providecommand{\Rat}{}\renewcommand{\Rat}{\operatorname{IsRat}}
\providecommand{\Res}{}\renewcommand{\Res}{\operatorname{Res}}
\providecommand{\Rlim}{}\renewcommand{\Rlim}{R\!\varprojlim}
\providecommand{\SIm}{}\renewcommand{\SIm}{\operatorname{SIm}}
\providecommand{\Sch}{}\renewcommand{\Sch}{\mathsf{Sch}}
\providecommand{\Sh}{}\renewcommand{\Sh}{\operatorname{Sh}}
\providecommand{\Span}{}\renewcommand{\Span}{\operatorname{span}}
\providecommand{\Spec}{}\renewcommand{\Spec}{\operatorname{Spec}}
\providecommand{\Tr}{}\renewcommand{\Tr}{\operatorname{Tr}}
\providecommand{\Trsm}{}\renewcommand{\Trsm}{\operatorname{Tr}^{\mathrm{sm}}}
\providecommand{\Wt}{}\renewcommand{\Wt}{\operatorname{Wt}}
\providecommand{\Z}{}\renewcommand{\Z}{\mathbb{Z}}
\providecommand{\Zar}{}\renewcommand{\Zar}{\mathrm{zar}}
\providecommand{\Zell}{}\renewcommand{\Zell}{{\mathbf{Z}_{\ell}}}
\providecommand{\Zl}{}\renewcommand{\Zl}{{\mathbf Z_\ell}}
\providecommand{\affet}{}\renewcommand{\affet}[1]{{#1}^{\mathrm{aff}}_{\mathrm{ét}}}
\providecommand{\affproet}{}\renewcommand{\affproet}[1]{{#1}^{\mathrm{aff}}_{\mathrm{proét}}}
\providecommand{\boxt}{}\renewcommand{\boxt}{\boxtimes_A^{\mathrm L}}
\providecommand{\cF}{}\renewcommand{\cF}{\mathcal{F}}
\providecommand{\cHom}{}\renewcommand{\cHom}{R\mathcal{H}\!om}
\providecommand{\cI}{}\renewcommand{\cI}{\mathcal{I}}
\providecommand{\cL}{}\renewcommand{\cL}{\mathcal L}
\providecommand{\cM}{}\renewcommand{\cM}{\mathcal M}
\providecommand{\cN}{}\renewcommand{\cN}{\mathcal N}
\providecommand{\cO}{}\renewcommand{\cO}{\mathcal O}
\providecommand{\cT}{}\renewcommand{\cT}{\mathcal T}
\providecommand{\checked}{}\renewcommand{\checked}{\color{red}\aftergroup\restorecolor}
\providecommand{\coeff}{}\renewcommand{\coeff}{\operatorname{coeff}}
\providecommand{\coker}{}\renewcommand{\coker}{\operatorname{coker}}
\providecommand{\colim}{}\renewcommand{\colim}{\operatorname{colim}}
\providecommand{\constructedby}{}\renewcommand{\constructedby}[1]{\par\smallskip\noindent\textcolor{purple!60!black}{\textsf{Constructed by:} }
  \footnotesize\sloppy\nolinkurl{#1}\par}
\providecommand{\deepresult}{}\renewcommand{\deepresult}[1]{\par\smallskip\noindent\textcolor{orange!45!black}{
  \textsf{Blueprint role: named deep theorem node; see the cluster proof architecture.}}
  \par\noindent\textcolor{orange!45!black}{\textsf{Mathematical source:} }#1\par}
\providecommand{\dt}{}\renewcommand{\dt}{\mathbin{\otimes^{\mathbf L}}}
\providecommand{\dw}{}\renewcommand{\dw}{\operatorname{dw}}
\providecommand{\et}{}\renewcommand{\et}[1]{{#1}_{\mathrm{ét}}}
\providecommand{\ev}{}\renewcommand{\ev}{\operatorname{ev}}
\providecommand{\fib}{}\renewcommand{\fib}{\operatorname{fib}}
\providecommand{\formalizationnote}{}\renewcommand{\formalizationnote}[1]{\begin{quote}\small\raggedright\textbf{Formalization note.} #1\end{quote}}
\providecommand{\fpqc}{}\renewcommand{\fpqc}{X_{\mathrm{fpqc}}}
\providecommand{\gammacont}{}\renewcommand{\gammacont}{\Gamma_{\mathrm{cont}}}
\providecommand{\gm}{}\renewcommand{\gm}{\mathbf G_m}
\providecommand{\hocolim}{}\renewcommand{\hocolim}{\operatorname{hocolim}}
\providecommand{\holim}{}\renewcommand{\holim}{\operatorname{holim}}
\providecommand{\hrB}{}\renewcommand{\hrB}{\@ifnextchar]{\hrEnd}{\hrStep}}
\providecommand{\hrEnd}{}\renewcommand{\hrEnd}{\@firstoftwo{}}
\providecommand{\hrStep}{}\renewcommand{\hrStep}{\@firstoftwo{\hrB}}
\providecommand{\id}{}\renewcommand{\id}{\operatorname{id}}
\providecommand{\kunFl}{}\renewcommand{\kunFl}{\mathbf F_\ell}
\providecommand{\kunZn}{}\renewcommand{\kunZn}{\mathbf Z/n\mathbf Z}
\providecommand{\kuncd}{}\renewcommand{\kuncd}{\operatorname{cd}}
\providecommand{\kundim}{}\renewcommand{\kundim}{\dim}
\providecommand{\kunotimes}{}\renewcommand{\kunotimes}[1]{\otimes^{\mathbf L}_{#1}}
\providecommand{\kunsupp}{}\renewcommand{\kunsupp}{\operatorname{Supp}}
\providecommand{\lotimes}{}\renewcommand{\lotimes}{\otimes^{\mathrm L}}
\providecommand{\m}{}\renewcommand{\m}{{\mathfrak{m}}}
\providecommand{\n}{}\renewcommand{\n}{{\mathfrak{n}}}
\providecommand{\nagataGap}{}\renewcommand{\nagataGap}[1]{\par\noindent\textbf{Formalization gap.} #1\par}
\providecommand{\openobligation}{}\renewcommand{\openobligation}[1]{\par\smallskip\noindent\textcolor{red!65!black}{\textsf{Open obligation:} }
  \footnotesize\sloppy\nolinkurl{#1}\par}
\providecommand{\p}{}\renewcommand{\p}{{\mathfrak{p}}}
\providecommand{\pr}{}\renewcommand{\pr}{\operatorname{pr}}
\providecommand{\preshv}{}\renewcommand{\preshv}[1]{\mathrm{PreShv}(#1)}
\providecommand{\pro}{}\renewcommand{\pro}[1]{\mathrm{Pro}(#1)}
\providecommand{\proet}{}\renewcommand{\proet}{X_{\mathrm{proét}}}
\providecommand{\q}{}\renewcommand{\q}{{\mathfrak{q}}}
\providecommand{\rad}{}\renewcommand{\rad}{\operatorname{rad}}
\providecommand{\rank}{}\renewcommand{\rank}{\operatorname{rank}}
\providecommand{\red}{}\renewcommand{\red}{\color{red}}
\providecommand{\restorecolor}{}\renewcommand{\restorecolor}{\color{black}}
\providecommand{\retiredobligation}{}\renewcommand{\retiredobligation}[1]{\par\smallskip\noindent\textcolor{red!75!black}{\textsf{Retired invalid route:} }
  \footnotesize\sloppy\nolinkurl{#1}\par}
\providecommand{\rk}{}\renewcommand{\rk}{\operatorname{rk}}
\providecommand{\sep}{}\renewcommand{\sep}{\textnormal{sep}}
\providecommand{\shv}{}\renewcommand{\shv}[1]{\mathrm{Shv}(#1)}
\providecommand{\spec}{}\renewcommand{\spec}[1]{\mathrm{Spec}(#1)}
\providecommand{\stacks}{}\renewcommand{\stacks}[1]{\href{https://stacks.math.columbia.edu/tag/#1}{\textnormal{Stacks, Tag~#1}}}
\providecommand{\stacksproject}{}\renewcommand{\stacksproject}[1]{{\cite[\href{https://stacks.math.columbia.edu/tag/#1}{Tag #1}]{stacks-project}}}
\providecommand{\statusnote}{}\renewcommand{\statusnote}[1]{\begin{quote}\small\raggedright\textbf{Status.} #1\end{quote}}
\providecommand{\tildeideal}{}\renewcommand{\tildeideal}[2]{\widetilde{(#1, #2)}}
\providecommand{\tildering}{}\renewcommand{\tildering}[3]{{#1}_{\widetilde{(#2, #3)}}}
\providecommand{\tr}{}\renewcommand{\tr}{\operatorname{tr}}
\providecommand{\wt}{}\renewcommand{\wt}{\operatorname{wt}}

\begin{document}
\pagestyle{plain}
\begin{center}
{\LARGE\bfseries Weil Conjectures\par}
\vspace{0.7em}
{\large Section 1: Introduction\par}
\vspace{0.5em}
{\small\color{bpmuted}Snapshot 2026-09-26, revision 8e4cb88e2a\par}
\end{center}
\vspace{0.5em}
{\small\tableofcontents}
\bigskip\hrule\bigskip

% object: chap-Introduction (blueprint/chap-Introduction/_section.md)
\section*{\texorpdfstring{1\quad Introduction}{1 Introduction}}
\phantomsection\label{bp1}%
\addcontentsline{toc}{section}{\texorpdfstring{1\quad Introduction}{1 Introduction}}
% text: prose-01-01
\phantomsection\label{bp2}%
These notes prove the Weil conjectures, which concern the number of solutions of polynomial equations over finite fields. Weil saw that rationality and the functional equation would follow from a cohomology theory for varieties over finite fields with a Lefschetz fixed point formula and Poincaré duality, while the deepest conjecture, an analogue of the Riemann hypothesis, needs in addition a bound on the eigenvalues of Frobenius. Following this plan, we construct such a theory, \(\ell\)-adic cohomology, and prove the trace formula and Poincaré duality for it. Then we prove the bound, which is the Deligne–Laumon theorem.

% object: sec-The-Weil-conjectures (blueprint/chap-Introduction/sec-The-Weil-conjectures/_section.md), content hidden
\subsection*{\texorpdfstring{1.1\quad The Weil conjectures}{1.1 The Weil conjectures}}
\phantomsection\label{bp3}%
\addcontentsline{toc}{subsection}{\texorpdfstring{1.1\quad The Weil conjectures}{1.1 The Weil conjectures}}
% text: prose-01-02
\phantomsection\label{bp4}%
Let \(X_0\) be a smooth projective geometrically connected curve of genus \(g\) over the finite field \(\mathbf F_q\), and let \(N_m\) be the number of its points with coordinates in \(\mathbf F_{q^m}\). The projective line has exactly \(1+q^m\) such points, and the \emph{Hasse–Weil bound} says that every curve of genus \(g\) comes close to this count:

\[
|N_m-(1+q^m)|\le 2g\,q^{m/2}.
\]

Hasse proved it for elliptic curves [\hyperref[bp21]{Has36}] and Weil for all curves [\hyperref[bp22]{Wei48a}].

To state Weil's conjectures in higher dimension we put all the numbers \(N_m\) into one generating function. For a separated scheme \(X_0\) of finite type over \(\mathbf F_q\), its \emph{zeta function} is

\[
Z(X_0,T)=\exp\Bigl(\sum_{m\ge 1}\#X_0(\mathbf F_{q^m})\frac{T^m}{m}\Bigr)
\in 1+T\mathbf Q[[T]].
\]

A closed point \(x\) of \(X_0\) with residue field \(\mathbf F_{q^{\deg x}}\) accounts for exactly \(\deg x\) points of \(X_0(\mathbf F_{q^m})\) when \(\deg x\) divides \(m\), and for none otherwise. Hence \(\#X_0(\mathbf F_{q^m})=\sum_{\deg x\mid m}\deg x\), so

\[
\sum_{m\ge1}\#X_0(\mathbf F_{q^m})\frac{T^m}{m}
=\sum_{x\in|X_0|}\sum_{k\ge1}\frac{T^{k\deg x}}{k}
=-\sum_{x\in|X_0|}\log\bigl(1-T^{\deg x}\bigr).
\]

Only finitely many closed points have a given degree, so exponentiating gives the Euler product

\[
Z(X_0,T)=\prod_{x\in|X_0|}\bigl(1-T^{\deg x}\bigr)^{-1},
\]

whose factors are power series with integer coefficients. So the coefficients of \(Z(X_0,T)\) are integers. For example, \(\#\mathbf P^1(\mathbf F_{q^m})=1+q^m\) gives \(Z(\mathbf P^1,T)=1/((1-T)(1-qT))\). Set \(T=q^{-s}\). Since the residue field of \(x\) has \(N(x)=q^{\deg x}\) elements, the Euler product becomes

\[
\zeta_{X_0}(s)=Z(X_0,q^{-s})=\prod_{x\in|X_0|}\bigl(1-N(x)^{-s}\bigr)^{-1},
\]

the analogue of Riemann's

\[
\zeta(s)=\sum_{n\ge1}n^{-s}=\prod_{p}\bigl(1-p^{-s}\bigr)^{-1},
\]

in which the closed points of \(\operatorname{Spec}\mathbf Z\) have residue fields with \(p\) elements. Riemann conjectured that the nontrivial zeros of \(\zeta(s)\) lie on the line \(\operatorname{Re}s=\tfrac12\). Weil conjectured the following.

\begin{enumerate}
\item \emph{Rationality.} \(Z(X_0,T)\) is a rational function of \(T\).
\item \emph{Functional equation.} If \(X_0\) is smooth and proper of pure dimension \(d\), then \(Z(X_0,1/(q^dT))=\pm q^{d\chi/2}T^{\chi}Z(X_0,T)\), where \(\chi=\sum_i(-1)^ib_i\) is the Euler characteristic, with \(b_i\) as in (4). For a curve of genus \(g\) it equals \(2-2g\).
\item \emph{Riemann hypothesis.} If \(X_0\) is smooth and proper of pure dimension \(d\), then \[
   Z(X_0,T)=\frac{P_1(T)P_3(T)\cdots P_{2d-1}(T)}
   {P_0(T)P_2(T)\cdots P_{2d}(T)},
   \] where \(P_i\in\mathbf Z[T]\) has constant term \(1\), and its reciprocal roots (the numbers \(\alpha\) with \(P_i(1/\alpha)=0\)) are algebraic integers all of whose complex conjugates have absolute value \(q^{i/2}\).
\item \emph{Betti numbers.} The degree \(b_i\) of \(P_i\) behaves like a Betti number. For instance, \(b_i=b_{2d-i}\), as for Poincaré duality on a compact manifold of real dimension \(2d\). More precisely, suppose that \(X_0\) comes from characteristic zero: there is a smooth projective variety \(Y\) over a number field \(K\) and a prime \(\mathfrak p\) of \(K\) with residue field \(\mathbf F_q\) such that \(Y\) extends to a smooth projective scheme over the local ring \(\mathcal O_{K,\mathfrak p}\) whose special fibre is \(X_0\). Then \(b_i\) is the \(i\)th Betti number of the compact complex manifold \(Y(\mathbf C)\), for any embedding \(K\subset\mathbf C\). For a curve of genus \(g\) this gives \(b_0=b_2=1\) and \(b_1=2g\).
\end{enumerate}

For a curve the conjectures recover the Hasse–Weil bound. Here \(Z(X_0,T)=P_1(T)/((1-T)(1-qT))\) with \(P_1\) of degree \(2g\), and statement (3) says that \(P_1(T)=\prod_{j=1}^{2g}(1-\alpha_jT)\) with \(|\alpha_j|=q^{1/2}\). Comparing the logarithmic derivatives of both sides with the definition of \(Z\) gives

\[
N_m=1+q^m-\sum_{j=1}^{2g}\alpha_j^m,
\]

and since each \(|\alpha_j^m|=q^{m/2}\), the Hasse–Weil bound follows. This is also why (3) is called the Riemann hypothesis. In the variable \(s\) with \(T=q^{-s}\), the zeros of \(P_1(q^{-s})\) satisfy \(|q^{-s}|=q^{-1/2}\), that is, they lie on the line \(\operatorname{Re}s=\tfrac12\), exactly where Riemann conjectured the zeros of \(\zeta(s)\) to lie.

We follow one example through the notes, the elliptic curve

\[
C_0\colon\ y^2=x^3-x\qquad\text{over }\mathbf F_3 .
\]

Counting points, including the point at infinity, gives \(N_1=4\), \(N_2=16\) and \(N_3=28\). Suppose the Weil conjectures hold for \(C_0\). Then \(Z(C_0,T)=P_1(T)/((1-T)(1-3T))\), where \(P_1\) has degree \(b_1=2g=2\). Its leading coefficient is \(q=3\), because by duality it is the determinant of Frobenius on \(H^1\) (Example \bpout{10.23}). So \(P_1=1-aT+qT^2\), and comparing the coefficients of \(T\) in \(Z\) gives \(a=1+q-N_1=0\). Hence \(P_1(T)=1+3T^2\), whose reciprocal roots \(\pm i\sqrt3\) have absolute value \(3^{1/2}\). Statements (1)–(4) alone would not fix the leading coefficient: \(1-3T^2\) also satisfies them, with the sign \(-\) in (2), and it is the count \(N_2=16\) that excludes it. As a check, the formula \(N_m=1+3^m-\bigl((i\sqrt3)^m+(-i\sqrt3)^m\bigr)\) gives \(16\) and \(28\) for \(m=2,3\). We return to this curve after duality, after the trace formula and at the end (Examples \bpout{10.23}, \bpout{Example 12.32} and \bpout{Example 19.9}).


% object: sec-Weil-cohomology (blueprint/chap-Introduction/sec-Weil-cohomology/_section.md), content hidden
\subsection*{\texorpdfstring{1.2\quad Weil cohomology}{1.2 Weil cohomology}}
\phantomsection\label{bp5}%
\addcontentsline{toc}{subsection}{\texorpdfstring{1.2\quad Weil cohomology}{1.2 Weil cohomology}}
% text: prose-01-03
\phantomsection\label{bp6}%
Write \(X=X_0\times_{\mathbf F_q}\overline{\mathbf F}_q\), and let \(F:X\to X\) be the \emph{Frobenius}: if \(X_0\) is cut out of an affine space by equations with coefficients in \(\mathbf F_q\), then \(F\) raises each coordinate to its \(q\)-th power. A point of \(X\) is fixed by \(F^m\) exactly when its coordinates lie in \(\mathbf F_{q^m}\), so \(\#X_0(\mathbf F_{q^m})\) is the number of fixed points of \(F^m\). Weil's idea was to count these fixed points as in topology.

Suppose that we had a cohomology theory with compact supports, assigning to \(X\) finite-dimensional vector spaces \(H^i_c(X)\) over a field of characteristic zero on which \(F\) acts. Compact supports are needed because the trace formula below has to hold for every \(X_0\), proper or not: point counts are additive when \(X_0\) is cut into an open subscheme and its complement, and compactly supported cohomology has the matching long exact sequence. For proper \(X_0\) we write \(H^i(X)=H^i_c(X)\); for the conjectures, which concern smooth proper varieties, this is the only cohomology we need. Suppose that the theory satisfies the following.

\begin{itemize}
\item (a) \emph{Finiteness.} \(H^i_c(X)=0\) unless \(0\le i\le 2\dim X\).
\item (b) \emph{Trace formula.} The fixed points of \(F^m\) are counted by a Lefschetz formula, \[
  \#X_0(\mathbf F_{q^m})=\sum_i(-1)^i\operatorname{Tr}\bigl(F^m\mid H^i_c(X)\bigr).
  \]
\item (c) \emph{Poincaré duality.} For smooth proper \(X_0\) of pure dimension \(d\), cup product gives a perfect pairing between \(H^i\) and \(H^{2d-i}\), and \(\langle Fx,Fy\rangle=q^d\langle x,y\rangle\). The factor \(q^d\) is the degree of \(F\). It is the scalar by which \(F\) acts on the one-dimensional space \(H^{2d}\) when \(X\) is connected.
\item (d) \emph{Purity.} For smooth proper \(X_0\), the eigenvalues of \(F\) on \(H^i(X)\) are algebraic integers all of whose complex conjugates have absolute value \(q^{i/2}\).
\end{itemize}

Then the conjectures follow. For rationality we compare logarithms. For an endomorphism \(F\) of a finite-dimensional vector space \(V\) we have

\[
\log\det(1-TF\mid V)^{-1}=\sum_{m\ge1}\operatorname{Tr}(F^m\mid V)\frac{T^m}{m},
\]

as one sees by triangulating \(F\): if \(F\) has eigenvalues \(\alpha_1,\dots,\alpha_b\), both sides equal \(\sum_j\sum_m\alpha_j^mT^m/m\). Substituting the trace formula (b) into the definition of \(Z\) therefore gives

\[
\log Z(X_0,T)=\sum_{m\ge1}\sum_i(-1)^i\operatorname{Tr}\bigl(F^m\mid H^i_c(X)\bigr)\frac{T^m}{m}
=\sum_i(-1)^{i+1}\log\det\bigl(1-TF\mid H^i_c(X)\bigr),
\]

and so

\[
Z(X_0,T)=\prod_i\det\bigl(1-TF\mid H^i_c(X)\bigr)^{(-1)^{i+1}},
\]

a finite product by (a). In particular \(Z(X_0,T)\) is rational, and for smooth proper \(X_0\) the polynomial \(P_i(T)\) of the conjectures is \(\det(1-TF\mid H^i(X))\). By duality, the eigenvalues in degree \(2d-i\) are the numbers \(q^d/\alpha\) for the eigenvalues \(\alpha\) in degree \(i\), and this gives the functional equation. Purity gives the Riemann hypothesis.

Purity also shows that the factors do not depend on the cohomology theory. Indeed, the factors in different degrees have reciprocal roots of different absolute values, so no cancellation can occur in the alternating product. Each factor is therefore determined by \(Z(X_0,T)\) alone, and so is its degree \(b_i\).

Which cohomology could this be? On an irreducible scheme a constant sheaf is flasque, so its Zariski cohomology vanishes in positive degrees. It cannot see the \(2g\) independent loops of a curve of genus \(g\). Coherent cohomology sees \(H^1(C,\mathcal O_C)\), of dimension \(g\). But this is a vector space over \(\overline{\mathbf F}_q\), and traces on it determine numbers of points only modulo \(p\).

An argument of Serre (see [\hyperref[bp23]{Den22}]) shows that the coefficient field cannot be \(\mathbf Q\), \(\mathbf R\) or \(\mathbf Q_p\) either. Let \(C\) be a supersingular elliptic curve over \(\overline{\mathbf F}_p\). Then \(D=\operatorname{End}(C)\otimes\mathbf Q\) is a quaternion algebra over \(\mathbf Q\), and \(D\otimes\mathbf R\) and \(D\otimes\mathbf Q_p\) are still division algebras [\hyperref[bp24]{Deu41}]. Endomorphisms of \(C\) would act on \(H^1(C)\), a vector space of dimension \(2g=2\) over the coefficient field \(K\). This gives a map of four-dimensional \(K\)-algebras \(D\otimes K\to M_2(K)\). The map is injective because \(D\otimes K\) is simple, hence it is an isomorphism. That is impossible if \(D\otimes K\) is a division algebra. For \(\ell\neq p\), on the other hand, \(D\otimes\mathbf Q_\ell\simeq M_2(\mathbf Q_\ell)\), and the argument does not exclude \(\mathbf Q_\ell\).

The theory that works is étale cohomology. Its idea is to replace the open subsets of the Zariski topology by étale maps \(U\to X\), which play the role of small open neighbourhoods in topology: like a covering space, an étale map is a local isomorphism in the sense of the implicit function theorem. A sheaf in this topology assigns a group to every étale map \(U\to X\), compatibly with restriction, and étale cohomology is the derived functor of global sections. Because there are many more étale maps than open subsets, a sheaf can now have nonzero cohomology in positive degrees.

For a curve we can already see the missing loops in degree one. Let \(C\) be a smooth projective connected curve of genus \(g\) over \(\overline{\mathbf F}_q\), and let \(n\) be prime to \(q\). Three étale sheaves are involved: \(\mathbf G_m\), which sends \(U\to C\) to the unit group \(\mathcal O(U)^\times\); the same sheaf with the \(n\)th power map \(t\mapsto t^n\); and its kernel \(\mu_n\), the sheaf of \(n\)th roots of unity. The sequence

\[
0\to\mu_n\to\mathbf G_m\xrightarrow{\;t\mapsto t^n\;}\mathbf G_m\to0
\]

is exact in the étale topology. Indeed, surjectivity of a map of sheaves only has to hold locally, and an \(n\)th root of a unit \(a\) exists on the étale cover defined by \(t^n=a\), which is étale because \(n\) is invertible. (In the Zariski topology the last map is not surjective.) As in the Zariski topology, \(H^1(C_{\acute{e}t},\mathbf G_m)\) is the Picard group of \(C\) (Hilbert's Theorem 90). The global units of \(C\) are \(\overline{\mathbf F}_q^\times\), in which \(n\)th roots exist, so \(H^0(\mathbf G_m)\xrightarrow{n}H^0(\mathbf G_m)\) is surjective. The long exact sequence therefore identifies \(H^1(C_{\acute{e}t},\mu_n)\) with the kernel of multiplication by \(n\) on \(H^1(C_{\acute{e}t},\mathbf G_m)=\operatorname{Pic}(C)\).

To compute this kernel we use the degree. It gives an exact sequence \(0\to\operatorname{Pic}^0(C)\to\operatorname{Pic}(C)\xrightarrow{\deg}\mathbf Z\to0\), and since \(\mathbf Z\) has no torsion, every line bundle killed by \(n\) has degree zero. So \(\operatorname{Pic}(C)[n]=\operatorname{Pic}^0(C)[n]\). The group \(\operatorname{Pic}^0(C)\) is the group of points of the \emph{Jacobian} \(J\) of \(C\), an abelian variety of dimension \(g\) (Section \bpout{C.4}), and on an abelian variety of dimension \(g\) the \(n\)-torsion points form a group \((\mathbf Z/n)^{2g}\) when \(n\) is prime to the characteristic (Proposition \bpout{B.14}). Altogether,

\[
H^1(C_{\acute{e}t},\mu_n)\simeq\operatorname{Pic}(C)[n]=J[n]\simeq(\mathbf Z/n)^{2g}
\]

(Proposition \bpout{D.2}), which is the size of \(H^1\) of a compact Riemann surface of genus \(g\) with coefficients \(\mathbf Z/n\). Now let \(n=\ell^k\) and let \(k\to\infty\). The groups \(J[\ell^k]\), with the maps \(J[\ell^{k+1}]\to J[\ell^k]\) given by multiplication by \(\ell\), form the \emph{Tate module} \(T_\ell J=\varprojlim_kJ[\ell^k]\), free of rank \(2g\) over \(\mathbf Z_\ell\) (Proposition \bpout{B.14}). On the cohomology side these maps correspond to \(\mu_{\ell^{k+1}}\to\mu_{\ell^k}\), \(z\mapsto z^\ell\), so the identification above passes to the limit:

\[
\varprojlim_kH^1(C_{\acute{e}t},\mu_{\ell^k})\simeq T_\ell J
\]

(Proposition \bpout{C.17}). In other words, the Tate module is the first cohomology of \(C\) with coefficients in the inverse system \((\mu_{\ell^k})_k\), which is written \(\mathbf Z_\ell(1)\) (we come back to what this limit means below). Weil used \(T_\ell J\) in his proof for curves, in place of an \(H^1\) that was not yet available.

Frobenius acts on both sides. Let \(\pi\) be the endomorphism of \(J\) induced by Frobenius; it raises the coordinates of torsion points to the \(q\)-th power, and under the isomorphism above it is the action of the arithmetic Frobenius \(\sigma:z\mapsto z^q\) in \(\operatorname{Gal}(\overline{\mathbf F}_q/\mathbf F_q)\) on the torsion points. The map \(F\) on cohomology is pullback along the Frobenius of \(C\), which is the inverse of this Galois action on \(\overline{\mathbf F}_q\)-points (we compare the two Frobenius maps in Section \bpout{12.1}). Hence \(F\) acts on \(\varprojlim_kH^1(C_{\acute{e}t},\mu_{\ell^k})\) by \(\pi^{-1}\). Weil showed [\hyperref[bp25]{Wei48b}] that the characteristic polynomial of \(\pi\) on \(T_\ell J\) has integer coefficients, and that the numerator of the zeta function is

\[
P_1(T)=\det(1-T\pi\mid T_\ell J).
\]

Because of the twist, \(F\) acts on the untwisted group \(H^1(C,\mathbf Q_\ell)\) by \(q\pi^{-1}\) (Section \bpout{10.5}). Its eigenvalues \(q/\alpha\) run through the eigenvalues \(\alpha\) of \(\pi\) again, by the functional equation, so also \(P_1(T)=\det(1-TF\mid H^1(C,\mathbf Q_\ell))\), as predicted in Section \hyperref[bp3]{1.1}. For the curve \(C_0\) above, \(\pi\) acts on \(T_\ell J\) with characteristic polynomial \(T^2+3\), and indeed \(P_1(T)=1+3T^2\).

The example suggests how to obtain a theory with coefficients of characteristic zero: compute étale cohomology with the finite coefficients \(\mathbf Z/\ell^k\), where the theory works, take the inverse limit over \(k\) to obtain \(\mathbf Z_\ell\)-modules, and finally invert \(\ell\). The limit has to be taken with care, as \(\operatorname{Spec}\mathbf F_q\) already shows. On it, étale sheaves are sets with a continuous action of the Galois group \(\widehat{\mathbf Z}\), and \(H^1\) with coefficients in an abelian group \(A\) is the group of continuous homomorphisms \(\widehat{\mathbf Z}\to A\). For \(A=\mathbf Z/\ell^k\) this is \(\mathbf Z/\ell^k\), and the limit of these groups is \(\mathbf Z_\ell\). But the étale topology sees \(\mathbf Z_\ell\) only as a discrete group, and a continuous homomorphism from the compact group \(\widehat{\mathbf Z}\) to the discrete torsion-free group \(\mathbf Z_\ell\) is zero, so the constant sheaf \(\mathbf Z_\ell\) has \(H^1=0\). What one wants, therefore, is the limit of the cohomology groups, not the cohomology of the limit sheaf. After inverting \(\ell\) this gives a theory with coefficients in \(\mathbf Q_\ell\). However, the eigenvalues of Frobenius lie in an algebraic closure \(\overline{\mathbf Q}_\ell\) of \(\mathbf Q_\ell\) rather than in \(\mathbf Q_\ell\) itself, and so we also allow coefficients in \(\overline{\mathbf Q}_\ell\), where characteristic polynomials split.

The resulting theory is \(\ell\)-adic cohomology, for a prime \(\ell\) not dividing \(q\). It was constructed by Grothendieck and his collaborators, who proved (a), (b) and (c) for it in SGA 4 and SGA 5 [\hyperref[bp26]{SGA4}], [\hyperref[bp27]{SGA5}]. Purity (d) is the hard part, and Deligne proved it [\hyperref[bp28]{Del74}], [\hyperref[bp29]{Del80}]. The general form of his theorem, on weights of \(Rf_!\), was later reproved by Laumon with the \(\ell\)-adic Fourier transform [\hyperref[bp30]{Lau87}] (see also [\hyperref[bp31]{Kat88}], [\hyperref[bp32]{KW01}]); we call it the Deligne–Laumon theorem. Our proof uses the Fourier transform as Laumon does, but it is organized differently in places.

Of statement (4) we prove that \(b_i=b_{2d-i}\). We also prove that \(b_i=\dim H^i(X,\mathbf Q_\ell)\) for every \(\ell\) not dividing \(q\); in particular this dimension does not depend on \(\ell\) (Theorem \bpout{19.7}). The comparison with the Betti numbers of a lift to characteristic zero needs the comparison theorem between étale and singular cohomology, which we do not prove.


% object: sec-The-function-sheaf-dictionary (blueprint/chap-Introduction/sec-The-function-sheaf-dictionary/_section.md), content hidden
\subsection*{\texorpdfstring{1.3\quad The function–sheaf dictionary}{1.3 The function–sheaf dictionary}}
\phantomsection\label{bp7}%
\addcontentsline{toc}{subsection}{\texorpdfstring{1.3\quad The function–sheaf dictionary}{1.3 The function–sheaf dictionary}}
% text: prose-01-04
\phantomsection\label{bp8}%
The proof of purity does not treat one variety at a time. It cuts a variety into pieces and fibres it over other varieties, and point counts behave simply under both operations. If \(Z_0\subset X_0\) is closed with complement \(U_0\), then \(\#X_0(\mathbf F_q)=\#U_0(\mathbf F_q)+\#Z_0(\mathbf F_q)\), and for a morphism \(f:X_0\to Y_0\),

\[
\#X_0(\mathbf F_q)=\sum_{y\in Y_0(\mathbf F_q)}\#f^{-1}(y)(\mathbf F_q).
\]

On the right we sum a function on \(Y_0\) that is not constant. To follow such sums with cohomology, we replace functions by sheaves.

The sheaves in question are \emph{constructible}: there is a decomposition of \(X\) into finitely many locally closed pieces on each of which the sheaf is locally constant with finite stalks. After passing to the limit, the stalks are finite-dimensional \(\mathbf Q_\ell\)-vector spaces instead. A \emph{Frobenius structure} on a constructible sheaf \(K\) on \(X\) is an isomorphism \(F^*K\simeq K\), and a sheaf with a Frobenius structure is a \emph{Weil sheaf}. The same definitions apply to complexes of sheaves (Section \bpout{10.5}). Let \(x\) be a point of \(X_0\) over \(\mathbf F_{q^m}\). The corresponding point \(\bar x\) of \(X\) is fixed by \(F^m\), so the stalk of \((F^m)^*K\) at \(\bar x\) is \(K_{\bar x}\). Iterating the Frobenius structure gives an isomorphism \((F^m)^*K\simeq K\), and on the stalk at \(\bar x\) it becomes an endomorphism \(\operatorname{Frob}_x\) of \(K_{\bar x}\). Its trace defines the \emph{trace function}

\[
t_K:X_0(\mathbf F_{q^m})\longrightarrow\mathbf Q_\ell,\qquad x\longmapsto
\operatorname{Tr}(\operatorname{Frob}_x\mid K_{\bar x}),
\]

where for a complex the trace is the alternating sum of the traces on the cohomology groups. The constant sheaf has trace function \(1\).

Operations on functions correspond to operations on sheaves. For example, pulling back a function along \(f\) corresponds to the inverse image \(f^*K\), and the product of two functions to the tensor product. Similarly, summing a function over the fibres of \(f\) corresponds to the direct image with compact supports \(Rf_!K\), whose stalk at a point \(\bar y\) is the compactly supported cohomology of \(K\) on the fibre over \(\bar y\). The resulting identity \(t_{Rf_!K}(y)=\sum_{f(x)=y}t_K(x)\) is a form of the trace formula (Corollary \bpout{12.26}). When \(f\) is the map to a point, it is (b) with coefficients in \(K\). The one further operation we need is a Fourier transform, explained in Section \hyperref[bp13]{1.6}. This correspondence between functions and sheaves is Grothendieck's \emph{function–sheaf dictionary} [\hyperref[bp33]{SGA4½}, Sommes trig.]:

\begin{center}\small
\begin{tabular}{ll}
\toprule
functions on \(X_0(\mathbf F_{q^m})\) & Weil complexes on \(X\) \\
\midrule
\(t\mapsto t\circ g\) & \(g^*K\) \\
\(t\cdot t'\) & \(K\otimes^{\mathbf L}L\) \\
\(t\mapsto\bigl(y\mapsto\sum_{f(x)=y}t(x)\bigr)\) & \(Rf_!K\) \\
\(t=t|_{U}+t|_{Z}\) & \(j_!j^*K\to K\to i_*i^*K\to\) \\
\(\sum_xt(x)\) & \(R\Gamma_c(X,K)\) with its Frobenius \\
\(t\mapsto\bigl(y\mapsto-\sum_xt(x)\psi(\operatorname{Tr}_{\mathbf F_{q^m}/\mathbf F_q}(xy))\bigr)\) & \(Rp_{2!}(p_1^*K\otimes\mathcal K_\psi)[1]\) \\
\bottomrule
\end{tabular}
\end{center}

Proposition \bpout{12.28} states and proves the first five lines. The last line is built from the first three, as we explain in Section \hyperref[bp13]{1.6}. Since \(Rf_!\) of a sheaf is a complex, we work in derived categories throughout.


% object: sec-The-six-operations (blueprint/chap-Introduction/sec-The-six-operations/_section.md), content hidden
\subsection*{\texorpdfstring{1.4\quad The six operations}{1.4 The six operations}}
\phantomsection\label{bp9}%
\addcontentsline{toc}{subsection}{\texorpdfstring{1.4\quad The six operations}{1.4 The six operations}}
% text: prose-01-05
\phantomsection\label{bp10}%
Duality has no simple counterpart for functions, and it needs two more operations: the internal Hom of sheaves and the \emph{exceptional inverse image} \(f^!\), a right adjoint of \(Rf_!\). With \(f^*\), \(Rf_*\), \(Rf_!\) and \(\otimes\) these are Grothendieck's \emph{six operations}. Part I of these notes, roughly their first third, constructs them for all the coefficient rings we need. To see why \(f^!\) is the right tool, consider Poincaré duality for a smooth variety \(X\) of pure dimension \(d\) over an algebraically closed field \(k\), with coefficients \(\Lambda=\mathbf Z/\ell^n\). It should say that cup product

\[
H_c^i(X,\Lambda)\times H^{2d-i}(X,\Lambda(d))\longrightarrow
H^{2d}_c(X,\Lambda(d))\longrightarrow\Lambda
\]

is a perfect pairing. Here \(\Lambda(d)=\mu_{\ell^n}^{\otimes d}\) is the \emph{Tate twist} (Definition \bpout{2.53}). The twist is needed because Frobenius does not act trivially on top cohomology. For instance, the map \(F:\mathbf P^1\to\mathbf P^1\) has degree \(q\). It pulls the class of a point back to \(q\) times itself, so \(F\) acts on \(H^2(\mathbf P^1,\Lambda)\) by \(q\). With our conventions \(F\) acts on \(\Lambda(1)\) by \(q^{-1}\) (Section \bpout{2.12}), and therefore it acts trivially on \(H^2(\mathbf P^1,\Lambda(1))\). In general the twist \((d)\) removes the factor \(q^d\) from the pairing.

For an oriented manifold \(V\) of dimension \(r\) there is a clean proof, due to Verdier. Let \(a:V\to\mathrm{pt}\) be the map to a point. The functor \(Ra_!\), cohomology with compact supports, has a right adjoint \(a^!\), and the adjunction gives, for every sheaf \(G\) on \(V\),

\[
R\!\operatorname{Hom}\bigl(R\Gamma_c(V,G),\Lambda\bigr)\simeq R\!\operatorname{Hom}(G,a^!\Lambda).
\]

This isomorphism comes from the adjunction alone and holds for every \(V\); what makes it a duality theorem is a computation of \(a^!\Lambda\). It is a complex of sheaves on \(V\), and its restriction to an open subset \(U\) is the corresponding complex for \(U\). So it can be computed in a chart \(\mathbf R^r\). There its only cohomology sheaf sits in degree \(-r\), with stalks \(\Hom(H^r_c(\mathbf R^r,\Lambda),\Lambda)\simeq\Lambda\). An orientation identifies these stalks consistently, and this gives \(a^!\Lambda\simeq\Lambda[r]\). For \(G=\Lambda\) the displayed isomorphism then reads \(R\!\operatorname{Hom}(R\Gamma_c(V,\Lambda),\Lambda)\simeq R\!\operatorname{Hom}(\Lambda,\Lambda[r])=R\Gamma(V,\Lambda)[r]\). Taking cohomology in degree \(-i\) and using that \(\Lambda\) is a field or, for \(\Lambda=\mathbf Z/\ell^n\), a self-injective ring, this says \(\Hom(H^i_c(V,\Lambda),\Lambda)\simeq H^{r-i}(V,\Lambda)\), which is Poincaré duality for \(V\).

We follow the same route. For the structure morphism \(a\) of a smooth variety \(X\) of dimension \(d\) we show \(a^!\Lambda\simeq\Lambda_X(d)[2d]\), again by a local computation (\emph{smooth purity}). Poincaré duality follows.

The main theorem of Part I is that \(f^!\) exists and carries constructible complexes to constructible complexes (Theorem \bpout{5.2}). Existence is a general fact. Indeed, to give a right adjoint of \(Rf_!\) is to find, for each \(N\), an object \(f^!N\) with \(\Hom(M,f^!N)=\Hom(Rf_!M,N)\), naturally in \(M\), and Brown's representability theorem provides such objects for an exact functor between triangulated categories that commutes with arbitrary direct sums. It needs one hypothesis on the source: it must be \emph{compactly generated}, that is, built by shifts, cones and direct sums from a set of objects \(M\) for which \(\Hom(M,-)\) commutes with direct sums. Such objects are called compact; in the derived category of a ring, for example, the perfect complexes are the compact objects. Constructible complexes are such objects, but they are not closed under direct sums. We therefore apply the theorem to the larger category built from them in this way.

It remains to show that \(f^!\) carries constructible complexes to constructible complexes, and this is geometric. Locally \(f\) is a closed immersion followed by a projection from affine space, and affine space is open in a product of projective lines, so two cases remain. For the projection \(\mathbf P^1_Y\to Y\), we compute \(f^!\) explicitly, whereas for a closed immersion \(i\) with open complement \(j\), a distinguished triangle computes \(i^!K\) from \(K\) and \(Rj_*j^*K\). This case therefore comes down to the theorem of Deligne that the ordinary direct image \(Rj_*\) preserves constructible complexes.

The six operations must also be available for \(\mathbf Z_\ell\)- and \(\mathbf Q_\ell\)-coefficients, compatibly with reduction modulo \(\ell^n\). This is a problem, because inverse limits behave badly on the étale site: on \(\operatorname{Spec}\mathbf F_q\), as we saw in Section \hyperref[bp5]{1.2}, the constant sheaf \(\mathbf Z_\ell\) has \(H^1=0\), although the limit of the groups \(H^1(\operatorname{Spec}\mathbf F_q,\mathbf Z/\ell^k)\) is \(\mathbf Z_\ell\). For this reason the classical theory works with compatible systems of sheaves modulo \(\ell^n\) and builds the derived categories by hand [\hyperref[bp34]{Eke90}]. We use instead the pro-étale site of Bhatt and Scholze [\hyperref[bp35]{BS15}]. Its objects over \(X\) include, besides the étale maps \(U\to X\), inverse limits \(\varprojlim_iU_i\to X\) of affine étale maps, for example the spectrum of an algebraic closure of \(\mathbf F_q\) over \(\operatorname{Spec}\mathbf F_q\), and its coverings are the families that are jointly surjective in a strong sense. Every étale sheaf defines a sheaf on it, and every bounded below complex of étale sheaves has the same cohomology on the two sites (Chapter \bpout{2}). The gain is that \(\mathbf Z_\ell\), with its topology, becomes a sheaf whose cohomology is the limit: on the pro-étale site \(H^1(\operatorname{Spec}\mathbf F_q,\mathbf Z_\ell)=\mathbf Z_\ell\) (Example \bpout{2.57}). On it, \(\mathbf Z_\ell\) is a sheaf of rings whose sections are continuous functions to \(\mathbf Z_\ell\), so that an \(\ell\)-adic sheaf is simply a sheaf of modules over it, and we use the usual derived category of these modules, with no inverse systems. In this setting all the geometric arguments are made with coefficients \(\mathbf Z/\ell^n\), where the classical theory applies, and the pro-étale site then carries them over to \(\mathbf Z_\ell\), \(\mathbf Q_\ell\) and \(\overline{\mathbf Q}_\ell\).


% object: sec-The-trace-formula (blueprint/chap-Introduction/sec-The-trace-formula/_section.md), content hidden
\subsection*{\texorpdfstring{1.5\quad The trace formula}{1.5 The trace formula}}
\phantomsection\label{bp11}%
\addcontentsline{toc}{subsection}{\texorpdfstring{1.5\quad The trace formula}{1.5 The trace formula}}
% text: prose-01-06
\phantomsection\label{bp12}%
Part II proves Poincaré duality and the trace formula. Why should the Lefschetz formula hold for Frobenius? For a self-map \(\varphi\) of a compact manifold \(V\) whose fixed points are transverse, the Lefschetz fixed point formula reads

\[
\sum_i(-1)^i\operatorname{Tr}\bigl(\varphi^*\mid H^i(V,\mathbf Q)\bigr)
=\sum_{\varphi(x)=x}\operatorname{sign}\det\bigl(1-d\varphi_x\bigr),
\]

so each fixed point \(x\) is counted with the sign of \(\det(1-d\varphi_x)\). In our case the differential of Frobenius is zero, because \(d(u^q)=qu^{q-1}du=0\) in characteristic \(p\). So every fixed point should count with \(+1\). On a curve one can see this directly. Let \(u\) be a local coordinate at a fixed point, and \(v\) the same coordinate on the second factor of \(X\times X\). In these coordinates the graph of \(F^m\) is \(v=u^{q^m}\), and it meets the diagonal \(v=u\) with multiplicity one.

The Grothendieck–Lefschetz trace formula states that for every Weil complex \(K\),

\[
\operatorname{Tr}\bigl(F^m\mid R\Gamma_c(X,K)\bigr)
=\sum_{y\in X_0(\mathbf F_{q^m})}t_K(y)
\]

(Theorem \bpout{12.24}). For the constant sheaf it is axiom (b). In Example \bpout{12.32} we apply it to the elliptic curve \(C_0\colon y^2=x^3-x\) over \(\mathbf F_3\) of Section \hyperref[bp3]{1.1}, where it reads \(N_m=1-\operatorname{Tr}(F^m\mid H^1)+3^m\). Since \(N_1=4\), the trace of \(F\) on \(H^1\) is \(0\). Duality gives \(\det(F\mid H^1)=3\), and we recover \(Z(C_0,T)=(1+3T^2)/((1-T)(1-3T))\).

We deduce the trace formula, as in SGA 5 [\hyperref[bp27]{SGA5}, III, III B, XV], from a general Lefschetz formula. A \emph{correspondence} on \(X\) is a pair of maps \(c_1,c_2:C\to X\). For instance, the graph of a map \(\varphi\) is a correspondence with \(c_1=\mathrm{id}\) and \(c_2=\varphi\). To act on the cohomology of \(K\), a correspondence needs a map \(u:c_2^*K\to c_1^!K\), and for the graph of \(F\) this is the Frobenius structure \(F^*K\to K\). For such data the Lefschetz–Verdier formula (Theorem \bpout{11.13}) holds when \(X\) and \(C\) are proper. It writes the trace on cohomology as a sum, over the fixed points, of \emph{local terms}, numbers computed from \(K\) and \(u\) near each fixed point. In general, however, a local term is not the trace on the stalk.

To apply this to Frobenius on a curve \(X\), which need not be proper, we compactify \(X\) by adding finitely many points and extend \(K\) by zero. The added points are again fixed by Frobenius, but there \(K\) is no longer locally constant: its stalk is zero, while nearby stalks need not be. The main computation of Part II shows that even so, at every fixed point \(y\) of Frobenius on a curve, the local term is the trace of \(\operatorname{Frob}_y\) on \(K_{\bar y}\) (Proposition \bpout{12.16}). To prove this, we cut \(K\) into its restrictions to \(y\) and to the complement of \(y\). Local terms are additive for this decomposition, and for the piece supported at \(y\) the local term is the trace. It therefore suffices to show that the local term vanishes when \(K_{\bar y}=0\).

Here \(dF=0\) is used once more. Because of it, at \((y,y)\) the graph \(v=u^{q^m}\) is tangent to the axis \(v=0\), that is, to \(X\times\{y\}\), while the diagonal \(v=u\) is not. Blowing up the point \((y,y)\) of \(X\times X\) therefore separates the two curves. After the blow-up the local term is a map between cohomology groups supported near two disjoint curves meeting the exceptional divisor, and such a map factors through the stalk of \(K\) at \(y\); so it vanishes when \(K_{\bar y}=0\) (Lemma \bpout{12.15}). (In Chapter \bpout{12} we use instead the transposed graph \(\{(F^mx,x)\}\), which is tangent to \(\{y\}\times X\).)

Curves suffice for the general case. Indeed, the trace formula is compatible with \(Rf_!\), and a variety can be fibred in curves over a variety of smaller dimension.


% object: sec-The-Deligne-Laumon-theorem (blueprint/chap-Introduction/sec-The-Deligne-Laumon-theorem/_section.md), content hidden
\subsection*{\texorpdfstring{1.6\quad The Deligne–Laumon theorem}{1.6 The Deligne–Laumon theorem}}
\phantomsection\label{bp13}%
\addcontentsline{toc}{subsection}{\texorpdfstring{1.6\quad The Deligne–Laumon theorem}{1.6 The Deligne–Laumon theorem}}
% text: prose-01-07
\phantomsection\label{bp14}%
Part III proves purity. Fix an embedding \(\iota:\overline{\mathbf Q}_\ell\to\mathbf C\), and write \(q_x\) for the number of elements of the residue field of a closed point \(x\) of \(X_0\). A Weil sheaf with coefficients in \(\overline{\mathbf Q}_\ell\) is \emph{pure of weight \(w\)} if every eigenvalue \(\alpha\) of Frobenius on every stalk satisfies \(|\iota(\alpha)|=q_x^{w/2}\), and \emph{mixed} if it has a finite filtration whose quotients are pure. Similarly, a complex has \emph{weights at most \(w\)} if each cohomology sheaf \(\mathcal H^i\) is mixed with weights at most \(w+i\) (Definition \bpout{13.1}). The shift by \(i\) is chosen so that both the constant sheaf, pure of weight \(0\), and its cohomology complex \(R\Gamma(X,\overline{\mathbf Q}_\ell)\), whose \(H^i\) should have weight \(i\), have weights at most \(0\).

The \emph{Deligne–Laumon theorem} (Theorem \bpout{13.11}), proved by Deligne in Weil II [\hyperref[bp29]{Del80}] and reproved by Laumon with the Fourier transform [\hyperref[bp30]{Lau87}], says that if \(K\) has weights at most \(w\), so does \(Rf_!K\) for every morphism \(f\) of separated \(\mathbf F_q\)-schemes of finite type. It gives only an upper bound. For the constant sheaf on a smooth proper \(X_0\), duality supplies the lower one. With \(f\) the map to a point, the theorem says that the eigenvalues on \(H^i\) satisfy \(|\iota(\alpha)|\le q^{i/2}\). Since Poincaré duality pairs \(\alpha\) on \(H^i\) with \(q^d/\alpha\) on \(H^{2d-i}\), the same bound in degree \(2d-i\) gives \(|\iota(\alpha)|\ge q^{i/2}\). The two bounds together are purity. Theorem \bpout{19.7} deduces the Weil conjectures from it.

The proof has four steps, each of which has a model in classical mathematics.

\textbf{Step 1: from varieties to curves.} The theorem is stated for all sheaves, not only the constant one, because the proof is an induction on dimension. Away from a closed subset, a variety is fibred in curves over a variety of smaller dimension, much as one computes the cohomology of a fibre bundle by the Leray spectral sequence from the cohomology of the fibres and of the base. Along such a fibration, however, the direct image of the constant sheaf is no longer constant, so the induction needs the theorem for all sheaves. It thereby reduces to curves, provided the direct images that occur along the way are mixed (Theorem \bpout{18.8}). On a curve the statement says: if \(\mathcal F\) is a lisse sheaf, pure of weight \(w\), on a smooth curve \(U_0\), then \(H^1_c(U,\mathcal F)\) has weights at most \(w+1\) (Theorem \bpout{17.15}). Since every curve, after removing finitely many points, maps to the affine line by a finite map, it is enough to treat open subsets of \(\mathbf A^1\). The proof on curves combines two ideas: Rankin's positivity argument, which Deligne used in Weil I and Weil II [\hyperref[bp28]{Del74}], [\hyperref[bp29]{Del80}], and the Fourier transform, which Laumon [\hyperref[bp30]{Lau87}] used in place of Deligne's monodromy arguments.

\textbf{Step 2: positivity bounds the eigenvalues at points.} The model is Rankin's method for the coefficients of modular forms [\hyperref[bp36]{Ran39}]: the Dirichlet series \(\sum|\tau(n)|^2n^{-s}\) has nonnegative coefficients, so the position of its first pole bounds each individual coefficient. Here the role of the Dirichlet series is played by an \(L\)-function. Let \(\mathcal G\) be a \emph{lisse} sheaf on an affine curve \(U_0\), that is, a locally constant one. Such a sheaf is the same as a representation of the fundamental group, and its \emph{monodromy} is the image of that representation. Call \(\mathcal G\) \emph{real} if the characteristic polynomials of Frobenius on its stalks have real coefficients (Definition \bpout{14.21}). Its \(L\)-function is the product, over closed points \(x\), of the \emph{Euler factors} \(\det(1-\operatorname{Frob}_x T^{\deg x}\mid\mathcal G_{\bar x})^{-1}\). On an affine curve \(H^0_c=0\) and \(H^i_c=0\) for \(i>2\), so the trace formula shows that the \(L\)-function equals \(\det(1-TF\mid H^1_c)/\det(1-TF\mid H^2_c)\). Two facts make this useful.

\begin{itemize}
\item Consider an even tensor power \(\mathcal G^{\otimes 2m}\). The logarithm of an Euler factor has coefficients \(\operatorname{Tr}(\operatorname{Frob}_x^n\mid\mathcal G_{\bar x})^{2m}/n\), and these are nonnegative because the traces are real. Each Euler factor therefore has nonnegative coefficients bounded by those of the whole product. In particular it converges wherever the product does (Lemma \bpout{14.24}).
\item The product has poles only where \(\det(1-TF\mid H^2_c)\) vanishes, and \(H^2_c\) is easy to compute: it is the largest quotient of a stalk on which the monodromy acts trivially, twisted by \((-1)\).
\end{itemize}

The second fact bounds the eigenvalues on \(H^2_c\) of \(\mathcal G^{\otimes2m}\), and the first turns this into a bound at every point. To state the result, let \(\beta\) be the largest \emph{determinant weight} of the irreducible constituents of \(\mathcal G\); the determinant weight of a constituent of rank \(r\) is the weight of its determinant divided by \(r\) (Definition \bpout{14.7}). The eigenvalues on \(H^2_c\) of \(\mathcal G^{\otimes2m}\) have absolute value at most \(q^{m\beta+1}\). Indeed, by the second fact \(H^2_c\) is a quotient of a stalk of \(\mathcal G^{\otimes2m}\) on which the geometric monodromy acts trivially, twisted by \((-1)\). On such a quotient Frobenius acts through one-dimensional pieces, each a constituent of \(\mathcal G^{\otimes2m}\), and the determinant weight of every constituent of \(\mathcal G^{\otimes2m}\) is at most \(2m\beta\) (Proposition \bpout{14.19}). So these eigenvalues have absolute value at most \(q^{m\beta}\), and the twist \((-1)\) multiplies them by \(q\). The product therefore converges for \(|T|<q^{-m\beta-1}\), and by the first fact the Euler factor at \(x\) converges there too. This factor has a pole where \(T^{\deg x}=\alpha^{-2m}\), for each eigenvalue \(\alpha\) of \(\operatorname{Frob}_x\) on \(\mathcal G_{\bar x}\). Hence \(|\iota(\alpha)|^{2m}\le q_x^{m\beta+1}\) for every \(m\), and taking \(2m\)-th roots and letting \(m\to\infty\) gives the intermediate result:

\begin{quote}
\emph{every Frobenius eigenvalue \(\alpha\) at a point of a real lisse sheaf satisfies \(|\iota(\alpha)|^2\le q_x^\beta\)}
\end{quote}

(Proposition \bpout{14.25}).

\textbf{Step 3: rank one, and purity of real sheaves.} The bound of Step 2 is only useful if the determinants have weights, that is, if lisse sheaves of rank one are pure. A rank-one sheaf is a character of the fundamental group, the analogue of a Hecke character, and it is pure as soon as its restriction to the geometric fundamental group has finite image. We show that this image is finite. If it were not, the logarithm of a power of the character would give a nonzero class in \(H^1\) of the curve fixed by Frobenius. This class cannot come from the removed points, since Frobenius acts on those classes by powers of \(q\). So it comes from the compactification, and it gives a nonzero vector in the Tate module of its Jacobian fixed by Frobenius. Such a vector is a compatible system of rational \(\ell\)-power torsion points of unbounded order, which is impossible because the Jacobian has only finitely many rational points (Section \bpout{14.1}). With the determinants under control, Step 2 bounds the eigenvalues of \(\mathcal G\) by its largest determinant weight. To bound those of a single constituent we apply Step 2 to exterior powers of \(\mathcal G\), in which that constituent is multiplied by the determinants of the others, which are now known; this isolates its weight (as in the example \(\overline{\mathbf Q}_\ell\oplus\overline{\mathbf Q}_\ell(-1)\) of Chapter \bpout{14}). Applied also to the dual, which gives the lower bound, it yields the second intermediate result:

\begin{quote}
\emph{each irreducible constituent of a real lisse sheaf on a curve is pure}
\end{quote}

(Theorem \bpout{14.26}).

\textbf{Step 4: the Fourier transform turns cohomology into a stalk.} Steps 2 and 3 bound eigenvalues at the points of a curve, but we must bound eigenvalues on its cohomology. The Fourier transform turns one into the other. Its model is the discrete Fourier transform \(\hat t(y)=\sum_{x\in\mathbf F_q}t(x)\psi(xy)\) of a function \(t\) on \(\mathbf F_q\), where \(\psi\) is a nontrivial additive character. Two properties matter. The value \(\hat t(0)\) is the sum of all values of \(t\), so a sum over the line becomes a single value. And Plancherel's formula \(\sum_y|\hat t(y)|^2=q\sum_x|t(x)|^2\) holds.

The classical example is the Gauss sum. Let \(\chi\) be a nontrivial multiplicative character, extended by \(\chi(0)=0\), and let \(g=\sum_x\chi(x)\psi(x)\) be its Gauss sum. Substituting \(x\mapsto x/y\) shows that \(\hat\chi(y)=\overline\chi(y)g\) for \(y\neq0\), so the transform has the same absolute value \(|g|\) at all \(q-1\) nonzero points, and \(\hat\chi(0)=0\). Plancherel now reads \((q-1)|g|^2=q(q-1)\), and we find \(|g|=\sqrt q\). In other words, a function of absolute value \(1\) is carried to one of absolute value \(\sqrt q\), that is, weight \(0\) to weight \(1\). This is the shape of the whole argument: first the transform is shown to have a single absolute value (to be pure), and then Plancherel finds which one.

Grothendieck's dictionary turns the transform into an operation on sheaves on the affine line. Read the formula for \(\hat t\) from the inside out: pull \(t\) back to the plane along \(p_1(x,y)=x\), multiply by the function \(\psi(xy)\), and sum over the fibres of \(p_2(x,y)=y\). Each step is a line of the table in Section \hyperref[bp7]{1.3}. So if \(\mathcal K_\psi\) is a sheaf on the plane with trace function \(\psi(xy)\) and we put

\[
\operatorname{FT}_\psi K=Rp_{2!}(p_1^*K\otimes\mathcal K_\psi)[1],
\]

then for \(y\in\mathbf F_q\) the table gives

\[
t_{\operatorname{FT}_\psi K}(y)
=-\,t_{Rp_{2!}(p_1^*K\otimes\mathcal K_\psi)}(y)
=-\sum_{x\in\mathbf F_q}t_{p_1^*K}(x,y)\,t_{\mathcal K_\psi}(x,y)
=-\sum_{x\in\mathbf F_q}t_K(x)\,\psi(xy)
=-\,\hat t_K(y).
\]

The sign comes from the shift \([1]\), which changes the sign of every trace.

Now let \(M\) be an irreducible lisse sheaf, pure of weight \(w\), on a dense open \(j:U_0\to\mathbf A^1\), and let \(j_*M\) be its direct image to \(\mathbf A^1\). Its compactly supported cohomology lives in degrees \(1\) and \(2\); as in Step 2, \(H^2_c\) is easy, and the group we must bound is \(H^1_c(\mathbf A^1,j_*M)\). Just as \(\hat t(0)\) is the sum of \(t\), the stalk at \(0\) of the Fourier transform of \(j_*M\) is the compactly supported cohomology of \(j_*M\). The sheaves \(\mathcal L_{\psi(ax)}\) with trace function \(x\mapsto\psi(ax)\) are exceptions, as the characters \(\psi(ax)\) are for the discrete transform: their transform is concentrated at one point. Suppose that no constituent of \(M\) over \(\overline{\mathbf F}_q\) is of this form. Then the transform is again of the form \(j'_*M'\) for a lisse sheaf \(M'\) on a dense open of the dual line (Proposition \bpout{16.17}). The hypothesis excludes in particular constant constituents (the case \(a=0\)), so the monodromy of \(M\) has no coinvariants and \(H^2_c(\mathbf A^1,j_*M)=0\). This gives the third intermediate result:

\begin{quote}
\emph{the stalk of \(j'_*M'\) at \(0\) is exactly \(H^1_c(\mathbf A^1,j_*M)\).}
\end{quote}

It remains to find the weight of \(M'\), and here the two halves of the Gauss sum argument return. To show that \(M'\) is pure we use Step 3: we twist \(M\) to weight \(0\) and add its dual, whose eigenvalues are the complex conjugates. The transform of the result is again a real sheaf, and it maps onto a twist of \(M'\), which is irreducible. So a twist of \(M'\) is an irreducible constituent of a real sheaf, hence pure by Step 3, and therefore \(M'\) is pure (Proposition \bpout{17.4}). The weight of \(M'\) is then found from the sheaf version of Plancherel's formula, which compares the sums of \(|t_M(x)|^2\) and \(|t_{M'}(y)|^2\) over \(\mathbf F_{q^m}\) as \(m\) grows: it is \(w+1\) (Theorem \bpout{17.5}). By Proposition \bpout{15.5}, the stalk of \(j'_*M'\) at \(0\) has weights at most \(w+1\), and this is the bound we want. For genus-one curves over \(\mathbf F_3\), the resulting purity gives \(1\le N_1\le7\) (Example \bpout{19.9}).


% object: sec-Outline (blueprint/chap-Introduction/sec-Outline/_section.md), content hidden
\subsection*{\texorpdfstring{1.7\quad Outline}{1.7 Outline}}
\phantomsection\label{bp15}%
\addcontentsline{toc}{subsection}{\texorpdfstring{1.7\quad Outline}{1.7 Outline}}
% text: prose-01-08
\phantomsection\label{bp16}%
The main theorem is the Weil conjectures for smooth proper varieties over \(\mathbf F_q\) (Theorem \bpout{19.7}), with the statement on Betti numbers as above. The proof follows Weil's plan: it constructs a cohomology theory satisfying the properties (a)–(d) of a Weil cohomology listed in Section \hyperref[bp5]{1.2}. Each part uses the one before it.

\begin{itemize}
\item Part I (Chapters \bpout{2}–\bpout{Chapter 8}) constructs the six operations for \(\ell\)-adic coefficients (Theorem \bpout{5.2}). It ends with smooth base change and the Künneth formula.
\item Part II (Chapters \bpout{9}–\bpout{Chapter 12}) proves smooth purity (Theorem \bpout{9.3}), Verdier and Poincaré duality (Theorem \bpout{10.10}), which is (c), and the trace formula, which is (b). It ends with rationality and the functional equation. As for (a), the groups \(H^i_c\) are finite because \(Rf_!\) preserves constructibility. They vanish outside \([0,2\dim X]\), and Theorem \bpout{10.15} treats the smooth proper case used for the Weil conjectures.
\item Part III (Chapters \bpout{13}–\bpout{Chapter 19}) defines weights, proves purity of real sheaves on curves, introduces the Fourier transform and proves the Deligne–Laumon theorem, first on curves and then in general. It ends with the Weil conjectures.
\end{itemize}

Most appendices prove results from algebraic geometry that the main text uses: Nagata compactification and Chow's lemma, abelian varieties, the Picard scheme of a curve, the cohomology of curves with torsion coefficients, and facts on algebraic groups. Appendix \bpout{E} is of a different kind. It is not a preliminary but a second, independent theory of local terms, by specialization to the normal cone, which computes the local terms of contracting correspondences; the main line of the proof does not use it.


% object: sec-Prerequisites-and-reading-guide (blueprint/chap-Introduction/sec-Prerequisites-and-reading-guide/_section.md), content hidden
\subsection*{\texorpdfstring{1.8\quad Prerequisites and reading guide}{1.8 Prerequisites and reading guide}}
\phantomsection\label{bp17}%
\addcontentsline{toc}{subsection}{\texorpdfstring{1.8\quad Prerequisites and reading guide}{1.8 Prerequisites and reading guide}}
% text: prose-01-09
\phantomsection\label{bp18}%
We assume a first course in algebraic geometry: schemes, quasi-coherent sheaves and their cohomology, curves and divisors, and étale morphisms; some homological algebra (derived functors, derived and triangulated categories); and singular cohomology with Poincaré duality and the Lefschetz fixed point formula, which serve as models. We do not assume étale cohomology. The results from the étale theory with torsion coefficients that we need [\hyperref[bp26]{SGA4}], [\hyperref[bp33]{SGA4½}] are stated where they are used, with references.

On a first reading one can take the foundations on trust and follow the main line:

\begin{itemize}
\item the statements of Chapter \bpout{3} and of the main theorem of Part I, Theorem \bpout{5.2};
\item the base change map and the statement of proper base change (Theorem \bpout{4.33});
\item Chapter \bpout{6} on why \(f^!\) preserves constructibility with torsion coefficients;
\item smooth purity with finite coefficients (Section \bpout{9.2}), which is where the twist \((d)\), the shift \([2d]\) and the trace map come from;
\item Chapter \bpout{10} on duality;
\item Frobenius and Weil complexes (Section \bpout{10.5}); in Chapter \bpout{12}, the local calculation at a fixed point of Frobenius on a curve;
\item Chapters \bpout{13}–\bpout{Chapter 19}, beginning with the plan in the opening of Part III.
\end{itemize}

The passages between coefficient rings (most of which are applications of Proposition \bpout{7.12}), the coherence of the various comparison maps, smooth base change (Chapter \bpout{8}), and the pro-étale foundations of Chapter \bpout{2} can be read when needed.


% object: sec-Conventions (blueprint/chap-Introduction/sec-Conventions/_section.md), content hidden
\subsection*{\texorpdfstring{1.9\quad Conventions}{1.9 Conventions}}
\phantomsection\label{bp19}%
\addcontentsline{toc}{subsection}{\texorpdfstring{1.9\quad Conventions}{1.9 Conventions}}
% text: prose-01-10
\phantomsection\label{bp20}%
We fix a prime \(\ell\) and a field \(k\) that is finite or algebraically closed, of characteristic different from \(\ell\). Unless said otherwise, schemes are separated and of finite type over \(k\), and morphisms are \(k\)-morphisms. The hypothesis on \(k\) is needed to bound the \(\ell\)-cohomological dimension of \(k\) (Lemma \bpout{2.46}), while the constructions of Sections \bpout{2.12} and \bpout{Section 3.1} make sense for arbitrary schemes. Finally, a citation such as [Stacks, \href{https://stacks.math.columbia.edu/tag/0F0C}{0F0C}] refers to the tag of a result in the Stacks Project [\hyperref[bp37]{Stacks}].

On the pro-étale site we regard every coefficient ring \(\Lambda\) as a condensed ring, that is, as a sheaf of rings on profinite sets, which remembers the topology of \(\Lambda\). (On the étale site, which we use only with finite coefficients, a finite ring is simply a ring.) The coefficient sheaf \(\Lambda_X\) on \(X\) is then the pullback of \(\Lambda\) along the map from \(X\) to a point (Section \bpout{2.12}). For a finite ring \(A\) this pullback is the classical constant sheaf.

Derived categories are the classical ones. Namely, \(D(X_{\mathrm{pro}\acute{e}t},\Lambda)\) is the derived category of the abelian category of sheaves of \(\Lambda_X\)-modules. On these categories all functors are derived, and we write them \(f^*\), \(Rf_*\), \(Rf_!\) and \(f^!\). Unless another category is named, we have \(f^*:D(Y_{\mathrm{pro}\acute{e}t},\Lambda)\to D(X_{\mathrm{pro}\acute{e}t},\Lambda)\) and \(Rf_*:D(X_{\mathrm{pro}\acute{e}t},\Lambda)\to D(Y_{\mathrm{pro}\acute{e}t},\Lambda)\), whereas the categories on which \(Rf_!\) and \(f^!\) act are fixed where these functors are defined. To avoid set-theoretic problems, we take sizes with respect to a fixed universe that is large enough to contain the w-contractible covers of Theorem \bpout{2.13}.

The table below lists the main notation, together with the technical terms (lisse, Weil structure, pure, mixed, real) whose meaning in these notes is fixed by a definition, and the place where each is defined. The Frobenius \(F\) of the Introduction, which raises coordinates to their \(q\)-th powers, is written \(\Phi_X\) in the body of the notes. A Frobenius structure \(F^*K\simeq K\) becomes there a Weil structure \(\operatorname{Fr}_X^*K\simeq K\) for the geometric Frobenius \(\operatorname{Fr}_X\), and Lemma \bpout{12.3} shows that the two give the same operators.

\textbf{Notation}

\begin{center}\small
\begin{tabularx}{\linewidth}{>{\raggedright\arraybackslash}X>{\raggedright\arraybackslash}X>{\raggedright\arraybackslash}X}
\toprule
symbol & meaning & defined in \\
\midrule
\(X_{\mathrm{pro}\acute{e}t}\) & the pro-étale site & \bpout{Definition 2.4} \\
\(\Lambda_X\), \(\Lambda_X(d)\) & coefficient sheaf, Tate twists & \bpout{Definition 2.53} \\
\(\Dbc(X,\Lambda)\), \(D^b_{\mathrm{ctf}}(X,\Lambda)\) & constructible complexes & \bpout{Definition 3.1} \\
lisse & locally of constant finite presentation & \bpout{Definition 3.2} \\
\(K_n\) & reduction modulo \(\pi_E^n\) & \bpout{Definition 2.61} \\
\(\operatorname{cd}_\ell\) & cohomological dimension & \bpout{Definition 2.45} \\
\(R\underline\Gamma_Z\) & local cohomology & \bpout{Definition 2.59} \\
\(\Loc(D^b_{\mathrm{ctf}}(X,\Lambda))\) & localizing subcategory & \bpout{Definition 5.1} \\
\(Rf_!\), \(f^!\) & exceptional operations & \bpout{Theorem 5.2} \\
\(\omega_f\) & relative dualizing complex & \bpout{Definition 9.4} \\
\(\operatorname{Tr}_f\) & trace map & \bpout{Definition 9.24} \\
\(\mathbb D_X\) & Verdier dual & \bpout{Definition 10.2} \\
\(\operatorname{Fr}_X\), Weil structure & geometric Frobenius, \(\operatorname{Fr}_X^*K\simeq K\) & \bpout{Definition 10.18} \\
\(\Phi_X\) & \(q\)-power Frobenius & \bpout{Definition 12.2} \\
\(c\), \(\Fix(c)\) & correspondences & \bpout{Definition 11.5} \\
\(T_c(u)\) & action of a correspondence & \bpout{Definition 11.11} \\
\(\operatorname{Tr}_c(u)\), \(\operatorname{LT}\) & trace class, local terms & \bpout{Definition 11.9} \\
pure, mixed & \(\iota\)-weights & \bpout{Definition 13.1} \\
\(\operatorname{dw}_\iota\) & determinant weight & \bpout{Definition 14.7} \\
real & real Weil sheaves & \bpout{Definition 14.21} \\
\(j_*\mathcal F\) & middle extension & \bpout{Definition 15.1} \\
\(\operatorname{FT}_\psi\) & Fourier transform & Chapter \bpout{16} \\
\(T_\ell A\), \(V_\ell A\) & Tate module & \bpout{Definition B.13} \\
\bottomrule
\end{tabularx}
\end{center}



\clearpage\phantomsection\hypertarget{bp-outside}{}%
\section*{Outside this export}
\addcontentsline{toc}{section}{Outside this export}
References marked \ensuremath{{}^\ast} are to parts of the blueprint this document leaves out:
\begin{description}
\item[Chapter 2] The pro-étale site, in Ⅰ The six operations (\bpslug{chap-{}\allowbreak{}The-{}\allowbreak{}pro-{}\allowbreak{}etale-{}\allowbreak{}site})
\item[Definition 2.4] The small pro-étale site, in Ⅰ The six operations (\bpslug{definition-{}\allowbreak{}small-{}\allowbreak{}pro-{}\allowbreak{}etale-{}\allowbreak{}site})
\item[Theorem 2.13] Existence of contractible covers, in Ⅰ The six operations (\bpslug{theorem-{}\allowbreak{}existence-{}\allowbreak{}contractible-{}\allowbreak{}covers})
\item[Definition 2.45] Cohomological dimension, in Ⅰ The six operations (\bpslug{definition-{}\allowbreak{}cohomological-{}\allowbreak{}dimension})
\item[Lemma 2.46] Cohomological dimension of the base field, in Ⅰ The six operations (\bpslug{lemma-{}\allowbreak{}cohomological-{}\allowbreak{}dimension-{}\allowbreak{}base-{}\allowbreak{}field})
\item[Section 2.12] Coefficient rings, in Ⅰ The six operations (\bpslug{sec-{}\allowbreak{}Coefficient-{}\allowbreak{}rings})
\item[Definition 2.53] Tate twists, in Ⅰ The six operations (\bpslug{definition-{}\allowbreak{}tate-{}\allowbreak{}twists})
\item[Example 2.57] The first cohomology of a finite field, in Ⅰ The six operations (\bpslug{example-{}\allowbreak{}first-{}\allowbreak{}cohomology-{}\allowbreak{}finite-{}\allowbreak{}field})
\item[Definition 2.59] Local cohomology, in Ⅰ The six operations (\bpslug{definition-{}\allowbreak{}local-{}\allowbreak{}cohomology})
\item[Definition 2.61] Derived adic completeness, in Ⅰ The six operations (\bpslug{definition-{}\allowbreak{}derived-{}\allowbreak{}adic-{}\allowbreak{}completeness})
\item[Chapter 3] Constructible complexes, in Ⅰ The six operations (\bpslug{chap-{}\allowbreak{}Constructible-{}\allowbreak{}complexes})
\item[Section 3.1] Derived functors, in Ⅰ The six operations (\bpslug{sec-{}\allowbreak{}Derived-{}\allowbreak{}functors})
\item[Definition 3.1] Constructible sheaves, in Ⅰ The six operations (\bpslug{definition-{}\allowbreak{}constructible-{}\allowbreak{}sheaves})
\item[Definition 3.2] Lisse sheaves, in Ⅰ The six operations (\bpslug{definition-{}\allowbreak{}lisse-{}\allowbreak{}sheaves})
\item[Theorem 4.33] Pro-étale proper base change, in Ⅰ The six operations (\bpslug{theorem-{}\allowbreak{}proper-{}\allowbreak{}base-{}\allowbreak{}change-{}\allowbreak{}classical-{}\allowbreak{}pro})
\item[Definition 5.1] The localizing subcategory generated by $\Dctf(X,\Lambda)$, in Ⅰ The six operations (\bpslug{definition-{}\allowbreak{}localizing-{}\allowbreak{}subcategory})
\item[Theorem 5.2] The exceptional inverse image, in Ⅰ The six operations (\bpslug{theorem-{}\allowbreak{}exceptional-{}\allowbreak{}inverse-{}\allowbreak{}image})
\item[Chapter 6] Constructibility of $f^!$, in Ⅰ The six operations (\bpslug{chap-{}\allowbreak{}Constructibility-{}\allowbreak{}of-{}\allowbreak{}f})
\item[Proposition 7.12] Isomorphisms reduce to finite coefficients, in Ⅰ The six operations (\bpslug{proposition-{}\allowbreak{}isomorphisms-{}\allowbreak{}reduce-{}\allowbreak{}finite-{}\allowbreak{}coefficients})
\item[Chapter 8] Smooth base change, in Ⅰ The six operations (\bpslug{chap-{}\allowbreak{}Smooth-{}\allowbreak{}base-{}\allowbreak{}change})
\item[Chapter 9] Smooth purity, in Ⅱ Duality and the trace formula (\bpslug{chap-{}\allowbreak{}Smooth-{}\allowbreak{}purity})
\item[Theorem 9.3] Smooth purity for all admissible coefficients, in Ⅱ Duality and the trace formula (\bpslug{theorem-{}\allowbreak{}smooth-{}\allowbreak{}purity-{}\allowbreak{}all-{}\allowbreak{}admissible-{}\allowbreak{}coefficients})
\item[Definition 9.4] Relative dualizing complex, in Ⅱ Duality and the trace formula (\bpslug{definition-{}\allowbreak{}relative-{}\allowbreak{}dualizing-{}\allowbreak{}complex})
\item[Section 9.2] The tensor product map, in Ⅱ Duality and the trace formula (\bpslug{sec-{}\allowbreak{}The-{}\allowbreak{}tensor-{}\allowbreak{}product-{}\allowbreak{}map})
\item[Definition 9.24] Trace map, in Ⅱ Duality and the trace formula (\bpslug{definition-{}\allowbreak{}trace-{}\allowbreak{}map})
\item[Chapter 10] Verdier and Poincaré duality, in Ⅱ Duality and the trace formula (\bpslug{chap-{}\allowbreak{}Verdier-{}\allowbreak{}and-{}\allowbreak{}Poincare-{}\allowbreak{}duality})
\item[Definition 10.2] Verdier dual and evaluation, in Ⅱ Duality and the trace formula (\bpslug{definition-{}\allowbreak{}verdier-{}\allowbreak{}dual-{}\allowbreak{}evaluation})
\item[Theorem 10.10] Global Verdier duality, in Ⅱ Duality and the trace formula (\bpslug{theorem-{}\allowbreak{}global-{}\allowbreak{}verdier-{}\allowbreak{}duality})
\item[Theorem 10.15] The trace isomorphism in top degree, in Ⅱ Duality and the trace formula (\bpslug{theorem-{}\allowbreak{}trace-{}\allowbreak{}isomorphism-{}\allowbreak{}top-{}\allowbreak{}degree})
\item[Section 10.5] Weil complexes, in Ⅱ Duality and the trace formula (\bpslug{sec-{}\allowbreak{}Weil-{}\allowbreak{}complexes})
\item[Definition 10.18] Weil structure, in Ⅱ Duality and the trace formula (\bpslug{definition-{}\allowbreak{}weil-{}\allowbreak{}structure})
\item[Example 10.23] The running example: duality, in Ⅱ Duality and the trace formula (\bpslug{example-{}\allowbreak{}running-{}\allowbreak{}example-{}\allowbreak{}duality})
\item[Definition 11.5] Correspondences, in Ⅱ Duality and the trace formula (\bpslug{definition-{}\allowbreak{}correspondence})
\item[Definition 11.9] Trace class and local terms, in Ⅱ Duality and the trace formula (\bpslug{definition-{}\allowbreak{}trace-{}\allowbreak{}class-{}\allowbreak{}local-{}\allowbreak{}terms})
\item[Definition 11.11] Action of a correspondence on cohomology, in Ⅱ Duality and the trace formula (\bpslug{definition-{}\allowbreak{}action-{}\allowbreak{}correspondence})
\item[Theorem 11.13] Proper pushforward and Lefschetz–Verdier, in Ⅱ Duality and the trace formula (\bpslug{theorem-{}\allowbreak{}proper-{}\allowbreak{}pushforward-{}\allowbreak{}lefschetz-{}\allowbreak{}verdier})
\item[Chapter 12] The trace formula, in Ⅱ Duality and the trace formula (\bpslug{chap-{}\allowbreak{}The-{}\allowbreak{}trace-{}\allowbreak{}formula})
\item[Section 12.1] Frobenius correspondences, in Ⅱ Duality and the trace formula (\bpslug{sec-{}\allowbreak{}Frobenius-{}\allowbreak{}correspondences})
\item[Definition 12.2] Frobenius correspondences and their operators, in Ⅱ Duality and the trace formula (\bpslug{definition-{}\allowbreak{}frobenius-{}\allowbreak{}correspondences-{}\allowbreak{}their-{}\allowbreak{}operators})
\item[Lemma 12.3] Scheme Frobenius and geometric Frobenius, in Ⅱ Duality and the trace formula (\bpslug{lemma-{}\allowbreak{}scheme-{}\allowbreak{}frobenius-{}\allowbreak{}geometric-{}\allowbreak{}frobenius})
\item[Lemma 12.15] Vanishing of the local term at a zero stalk, in Ⅱ Duality and the trace formula (\bpslug{lemma-{}\allowbreak{}vanishing-{}\allowbreak{}local-{}\allowbreak{}term-{}\allowbreak{}zero-{}\allowbreak{}stalk})
\item[Proposition 12.16] The Frobenius local term on a curve, in Ⅱ Duality and the trace formula (\bpslug{proposition-{}\allowbreak{}frobenius-{}\allowbreak{}local-{}\allowbreak{}term-{}\allowbreak{}curve})
\item[Theorem 12.24] Grothendieck–Lefschetz trace formula, in Ⅱ Duality and the trace formula (\bpslug{theorem-{}\allowbreak{}grothendieck-{}\allowbreak{}lefschetz-{}\allowbreak{}trace-{}\allowbreak{}formula})
\item[Corollary 12.26] Trace functions and direct image with compact support, in Ⅱ Duality and the trace formula (\bpslug{corollary-{}\allowbreak{}trace-{}\allowbreak{}functions-{}\allowbreak{}direct-{}\allowbreak{}image-{}\allowbreak{}compact})
\item[Proposition 12.28] The function–sheaf dictionary, in Ⅱ Duality and the trace formula (\bpslug{proposition-{}\allowbreak{}function-{}\allowbreak{}sheaf-{}\allowbreak{}dictionary})
\item[Example 12.32] The running example: the trace formula, in Ⅱ Duality and the trace formula (\bpslug{example-{}\allowbreak{}running-{}\allowbreak{}example-{}\allowbreak{}trace-{}\allowbreak{}formula})
\item[Chapter 13] Weights, in Ⅲ Weights and the Riemann hypothesis (\bpslug{chap-{}\allowbreak{}Weights})
\item[Definition 13.1] Weights, purity and mixedness, in Ⅲ Weights and the Riemann hypothesis (\bpslug{definition-{}\allowbreak{}weights})
\item[Theorem 13.11] Weights of $Rf_!$, in Ⅲ Weights and the Riemann hypothesis (\bpslug{theorem-{}\allowbreak{}weights-{}\allowbreak{}rf-{}\allowbreak{}shriek})
\item[Chapter 14] Purity of real sheaves, in Ⅲ Weights and the Riemann hypothesis (\bpslug{chap-{}\allowbreak{}Purity-{}\allowbreak{}of-{}\allowbreak{}real-{}\allowbreak{}sheaves})
\item[Section 14.1] Sheaves of rank one, in Ⅲ Weights and the Riemann hypothesis (\bpslug{sec-{}\allowbreak{}Sheaves-{}\allowbreak{}of-{}\allowbreak{}rank-{}\allowbreak{}one})
\item[Definition 14.7] Determinant weight, in Ⅲ Weights and the Riemann hypothesis (\bpslug{definition-{}\allowbreak{}determinant-{}\allowbreak{}weight})
\item[Proposition 14.19] Determinant weights of tensor operations, in Ⅲ Weights and the Riemann hypothesis (\bpslug{proposition-{}\allowbreak{}determinant-{}\allowbreak{}weights-{}\allowbreak{}tensor-{}\allowbreak{}operations})
\item[Definition 14.21] Real sheaves, in Ⅲ Weights and the Riemann hypothesis (\bpslug{definition-{}\allowbreak{}real-{}\allowbreak{}sheaves})
\item[Lemma 14.24] Positivity and convergence, in Ⅲ Weights and the Riemann hypothesis (\bpslug{lemma-{}\allowbreak{}positivity-{}\allowbreak{}convergence})
\item[Proposition 14.25] A bound by the largest determinant weight, in Ⅲ Weights and the Riemann hypothesis (\bpslug{proposition-{}\allowbreak{}bound-{}\allowbreak{}largest-{}\allowbreak{}determinant-{}\allowbreak{}weight})
\item[Theorem 14.26] Purity of real sheaves on curves, in Ⅲ Weights and the Riemann hypothesis (\bpslug{theorem-{}\allowbreak{}purity-{}\allowbreak{}real-{}\allowbreak{}sheaves-{}\allowbreak{}curves})
\item[Definition 15.1] Middle extension, in Ⅲ Weights and the Riemann hypothesis (\bpslug{definition-{}\allowbreak{}middle-{}\allowbreak{}extension})
\item[Proposition 15.5] Weights of $j_*$ at the boundary points, in Ⅲ Weights and the Riemann hypothesis (\bpslug{proposition-{}\allowbreak{}weights-{}\allowbreak{}j-{}\allowbreak{}boundary-{}\allowbreak{}points})
\item[Chapter 16] The Fourier transform, in Ⅲ Weights and the Riemann hypothesis (\bpslug{chap-{}\allowbreak{}The-{}\allowbreak{}Fourier-{}\allowbreak{}transform})
\item[Proposition 16.17] The Fourier transform of a middle extension, in Ⅲ Weights and the Riemann hypothesis (\bpslug{proposition-{}\allowbreak{}fourier-{}\allowbreak{}transform-{}\allowbreak{}middle-{}\allowbreak{}extension})
\item[Proposition 17.4] Purity of the lisse part of the Fourier transform, in Ⅲ Weights and the Riemann hypothesis (\bpslug{proposition-{}\allowbreak{}purity-{}\allowbreak{}lisse-{}\allowbreak{}part-{}\allowbreak{}fourier-{}\allowbreak{}transform})
\item[Theorem 17.5] The weight of the Fourier transform, in Ⅲ Weights and the Riemann hypothesis (\bpslug{theorem-{}\allowbreak{}weight-{}\allowbreak{}fourier-{}\allowbreak{}transform})
\item[Theorem 17.15] The weight bound on a curve, in Ⅲ Weights and the Riemann hypothesis (\bpslug{theorem-{}\allowbreak{}weight-{}\allowbreak{}bound-{}\allowbreak{}curve})
\item[Theorem 18.8] Real constructible sheaves are mixed, in Ⅲ Weights and the Riemann hypothesis (\bpslug{theorem-{}\allowbreak{}real-{}\allowbreak{}constructible-{}\allowbreak{}sheaves-{}\allowbreak{}are-{}\allowbreak{}mixed})
\item[Chapter 19] Proof of the Weil conjectures, in Ⅲ Weights and the Riemann hypothesis (\bpslug{chap-{}\allowbreak{}Proof-{}\allowbreak{}of-{}\allowbreak{}the-{}\allowbreak{}Weil-{}\allowbreak{}conjectures})
\item[Theorem 19.7] The Weil conjectures for smooth proper schemes, in Ⅲ Weights and the Riemann hypothesis (\bpslug{theorem-{}\allowbreak{}weil-{}\allowbreak{}conjectures-{}\allowbreak{}smooth-{}\allowbreak{}proper-{}\allowbreak{}schemes})
\item[Example 19.9] The running example: purity, in Ⅲ Weights and the Riemann hypothesis (\bpslug{example-{}\allowbreak{}running-{}\allowbreak{}example-{}\allowbreak{}purity})
\item[Definition B.13] Tate module, in Appendices (\bpslug{definition-{}\allowbreak{}tate-{}\allowbreak{}module})
\item[Proposition B.14] The integral Tate module, in Appendices (\bpslug{proposition-{}\allowbreak{}integral-{}\allowbreak{}tate-{}\allowbreak{}module})
\item[Section C.4] The Jacobian, in Appendices (\bpslug{sec-{}\allowbreak{}The-{}\allowbreak{}Jacobian})
\item[Proposition C.17] $H^1$ of a curve and the Tate module of its Jacobian, in Appendices (\bpslug{proposition-{}\allowbreak{}h-{}\allowbreak{}1-{}\allowbreak{}curve-{}\allowbreak{}tate-{}\allowbreak{}module})
\item[Proposition D.2] Smooth proper curves, torsion prime to the characteristic, in Appendices (\bpslug{proposition-{}\allowbreak{}smooth-{}\allowbreak{}proper-{}\allowbreak{}curves-{}\allowbreak{}torsion-{}\allowbreak{}prime})
\item[Appendix E] Specialization and contracting correspondences, in Appendices (\bpslug{chap-{}\allowbreak{}Specialization-{}\allowbreak{}and-{}\allowbreak{}contracting-{}\allowbreak{}correspondences})
\end{description}
\clearpage\phantomsection%
\section*{References}
\addcontentsline{toc}{section}{References}
\bpbibitem{BBD82}{bp38}{A. Beilinson, J. Bernstein, and P. Deligne. \emph{Faisceaux pervers}, volume 100 of \emph{Astérisque}. 1982.}

\bpbibitem{BS15}{bp35}{B. Bhatt and P. Scholze. The pro-étale topology for schemes. \emph{Astérisque}, 369:99–201, 2015. \href{https://people.mpim-bonn.mpg.de/scholze/proetale.pdf}{Author's manuscript}.}

\bpbibitem{Con07}{bp39}{B. Conrad. Deligne's notes on Nagata compactifications. \emph{Journal of the Ramanujan Mathematical Society}, 22:205–257, 2007. \href{https://math.stanford.edu/\~conrad/papers/nagatafinal.pdf}{Author's manuscript}.}

\bpbibitem{Del74}{bp28}{P. Deligne. La conjecture de Weil. I. \emph{Publ. Math. IHÉS}, 43:273–307, 1974.}

\bpbibitem{Del80}{bp29}{P. Deligne. La conjecture de Weil. II. \emph{Publ. Math. IHÉS}, 52:137–252, 1980.}

\bpbibitem{Den22}{bp23}{C. Deninger. There is no “Weil-”cohomology theory with real coefficients for arithmetic curves. \emph{Ann. Sc. Norm. Super. Pisa Cl. Sci.}, 2022. To appear. arXiv:\href{https://arxiv.org/abs/2204.02714}{2204.02714}.}

\bpbibitem{Deu41}{bp24}{M. Deuring. Die Typen der Multiplikatorenringe elliptischer Funktionenkörper. \emph{Abh. Math. Sem. Univ. Hamburg}, 14:197–272, 1941.}

\bpbibitem{Eke90}{bp34}{T. Ekedahl. On the adic formalism. In \emph{The Grothendieck Festschrift, Vol. II}, volume 87 of \emph{Progress in Mathematics}, pages 197–218. Birkhäuser, 1990.}

\bpbibitem{Has36}{bp21}{H. Hasse. Zur Theorie der abstrakten elliptischen Funktionenkörper. I, II, III. \emph{J. Reine Angew. Math.}, 175:55–62, 69–88, 193–208, 1936.}

\bpbibitem{Hum72}{bp40}{J. E. Humphreys. \emph{Introduction to Lie Algebras and Representation Theory}, volume 9 of \emph{Graduate Texts in Mathematics}. Springer, 1972.}

\bpbibitem{Ill94}{bp41}{L. Illusie. Autour du théorème de monodromie locale. \emph{Astérisque}, 223:9–57, 1994. §4. \href{https://www.numdam.org/item/AST_1994__223__9_0/}{Original article}.}

\bpbibitem{Kat88}{bp31}{N. M. Katz. Travaux de Laumon. \emph{Astérisque}, 161–162:105–132, 1988. Séminaire Bourbaki, exp. 691.}

\bpbibitem{KW01}{bp32}{R. Kiehl and R. Weissauer. \emph{Weil conjectures, perverse sheaves and \(\ell\)-adic Fourier transform}, volume 42 of \emph{Ergebnisse der Mathematik}. Springer, 2001.}

\bpbibitem{Lau87}{bp30}{G. Laumon. Transformation de Fourier, constantes d'équations fonctionnelles et conjecture de Weil. \emph{Publ. Math. IHÉS}, 65:131–210, 1987.}

\bpbibitem{Mil86a}{bp42}{J. S. Milne. \href{https://www.jmilne.org/math/xnotes/AVs.pdf}{Abelian Varieties}. In \emph{Arithmetic Geometry}. 1986. Corrected online version.}

\bpbibitem{Mil86b}{bp43}{J. S. Milne. \href{https://www.jmilne.org/math/xnotes/JVs.pdf}{Jacobian Varieties}. In \emph{Arithmetic Geometry}. 1986. Corrected online version.}

\bpbibitem{Ran39}{bp36}{R. A. Rankin. Contributions to the theory of Ramanujan's function \(\tau(n)\) and similar arithmetical functions. II. \emph{Proc. Cambridge Philos. Soc.}, 35:357–372, 1939.}

\bpbibitem{SGA4}{bp26}{M. Artin, A. Grothendieck, and J.-L. Verdier, editors. \emph{Théorie des topos et cohomologie étale des schémas}, volumes 269, 270, 305 of \emph{Lecture Notes in Mathematics}. Springer, 1972–1973. SGA 4.}

\bpbibitem{SGA4½}{bp33}{P. Deligne. \emph{Cohomologie étale}, volume 569 of \emph{Lecture Notes in Mathematics}. Springer, 1977. SGA 4½; \emph{Théorèmes de finitude}, §3, and \emph{Rapport sur la formule des traces}, §3, and \emph{Sommes trigonométriques}, §1.}

\bpbibitem{SGA5}{bp27}{A. Grothendieck, editor. \emph{Cohomologie \(\ell\)-adique et fonctions \(L\)}, volume 589 of \emph{Lecture Notes in Mathematics}. Springer, 1977. SGA 5.}

\bpbibitem{SGA7II}{bp44}{P. Deligne and N. Katz, editors. \emph{Groupes de monodromie en géométrie algébrique}, volume 340 of \emph{Lecture Notes in Mathematics}. Springer, 1973. SGA 7 II; Exposé XIII, (2.1.7).}

\bpbibitem{Stacks}{bp37}{The Stacks Project Authors. \href{https://stacks.math.columbia.edu}{The Stacks Project}. Individual results are cited by their stable tags in the text.}

\bpbibitem{Var05}{bp45}{Y. Varshavsky. Lefschetz–Verdier trace formula and a generalization of a theorem of Fujiwara. 2005. Version 2. arXiv:\href{https://arxiv.org/abs/math/0505564}{math/0505564}. \href{https://arxiv.org/abs/math/0505564v2}{Author's preprint}.}

\bpbibitem{Wei48a}{bp22}{A. Weil. \emph{Sur les courbes algébriques et les variétés qui s'en déduisent}. Hermann, Paris, 1948.}

\bpbibitem{Wei48b}{bp25}{A. Weil. \emph{Variétés abéliennes et courbes algébriques}. Hermann, Paris, 1948.}

\bpbibitem{Wei49}{bp46}{A. Weil. Numbers of solutions of equations in finite fields. \emph{Bull. Amer. Math. Soc.}, 55:497–508, 1949.}

\end{document}
